% Options for packages loaded elsewhere
% Options for packages loaded elsewhere
\PassOptionsToPackage{unicode}{hyperref}
\PassOptionsToPackage{hyphens}{url}
%
\documentclass[
  indonesian,
  12pt,
  a4paper,
]{report}
\usepackage{xcolor}
\usepackage[top=4cm,bottom=3cm,left=4cm,right=3cm,headheight=15pt]{geometry}
\usepackage{amsmath,amssymb}
\setcounter{secnumdepth}{3}
\usepackage{iftex}
\ifPDFTeX
  \usepackage[T1]{fontenc}
  \usepackage[utf8]{inputenc}
  \usepackage{textcomp} % provide euro and other symbols
\else % if luatex or xetex
  \usepackage{unicode-math} % this also loads fontspec
  \defaultfontfeatures{Scale=MatchLowercase}
  \defaultfontfeatures[\rmfamily]{Ligatures=TeX,Scale=1}
\fi
\usepackage{lmodern}
\ifPDFTeX\else
  % xetex/luatex font selection
  \setmainfont[]{TeX Gyre Termes}
\fi
% Use upquote if available, for straight quotes in verbatim environments
\IfFileExists{upquote.sty}{\usepackage{upquote}}{}
\IfFileExists{microtype.sty}{% use microtype if available
  \usepackage[]{microtype}
  \UseMicrotypeSet[protrusion]{basicmath} % disable protrusion for tt fonts
}{}
\usepackage{setspace}
% Make \paragraph and \subparagraph free-standing
\makeatletter
\ifx\paragraph\undefined\else
  \let\oldparagraph\paragraph
  \renewcommand{\paragraph}{
    \@ifstar
      \xxxParagraphStar
      \xxxParagraphNoStar
  }
  \newcommand{\xxxParagraphStar}[1]{\oldparagraph*{#1}\mbox{}}
  \newcommand{\xxxParagraphNoStar}[1]{\oldparagraph{#1}\mbox{}}
\fi
\ifx\subparagraph\undefined\else
  \let\oldsubparagraph\subparagraph
  \renewcommand{\subparagraph}{
    \@ifstar
      \xxxSubParagraphStar
      \xxxSubParagraphNoStar
  }
  \newcommand{\xxxSubParagraphStar}[1]{\oldsubparagraph*{#1}\mbox{}}
  \newcommand{\xxxSubParagraphNoStar}[1]{\oldsubparagraph{#1}\mbox{}}
\fi
\makeatother

\usepackage{color}
\usepackage{fancyvrb}
\newcommand{\VerbBar}{|}
\newcommand{\VERB}{\Verb[commandchars=\\\{\}]}
\DefineVerbatimEnvironment{Highlighting}{Verbatim}{commandchars=\\\{\}}
% Add ',fontsize=\small' for more characters per line
\usepackage{framed}
\definecolor{shadecolor}{RGB}{241,243,245}
\newenvironment{Shaded}{\begin{snugshade}}{\end{snugshade}}
\newcommand{\AlertTok}[1]{\textcolor[rgb]{0.68,0.00,0.00}{#1}}
\newcommand{\AnnotationTok}[1]{\textcolor[rgb]{0.37,0.37,0.37}{#1}}
\newcommand{\AttributeTok}[1]{\textcolor[rgb]{0.40,0.45,0.13}{#1}}
\newcommand{\BaseNTok}[1]{\textcolor[rgb]{0.68,0.00,0.00}{#1}}
\newcommand{\BuiltInTok}[1]{\textcolor[rgb]{0.00,0.23,0.31}{#1}}
\newcommand{\CharTok}[1]{\textcolor[rgb]{0.13,0.47,0.30}{#1}}
\newcommand{\CommentTok}[1]{\textcolor[rgb]{0.37,0.37,0.37}{#1}}
\newcommand{\CommentVarTok}[1]{\textcolor[rgb]{0.37,0.37,0.37}{\textit{#1}}}
\newcommand{\ConstantTok}[1]{\textcolor[rgb]{0.56,0.35,0.01}{#1}}
\newcommand{\ControlFlowTok}[1]{\textcolor[rgb]{0.00,0.23,0.31}{\textbf{#1}}}
\newcommand{\DataTypeTok}[1]{\textcolor[rgb]{0.68,0.00,0.00}{#1}}
\newcommand{\DecValTok}[1]{\textcolor[rgb]{0.68,0.00,0.00}{#1}}
\newcommand{\DocumentationTok}[1]{\textcolor[rgb]{0.37,0.37,0.37}{\textit{#1}}}
\newcommand{\ErrorTok}[1]{\textcolor[rgb]{0.68,0.00,0.00}{#1}}
\newcommand{\ExtensionTok}[1]{\textcolor[rgb]{0.00,0.23,0.31}{#1}}
\newcommand{\FloatTok}[1]{\textcolor[rgb]{0.68,0.00,0.00}{#1}}
\newcommand{\FunctionTok}[1]{\textcolor[rgb]{0.28,0.35,0.67}{#1}}
\newcommand{\ImportTok}[1]{\textcolor[rgb]{0.00,0.46,0.62}{#1}}
\newcommand{\InformationTok}[1]{\textcolor[rgb]{0.37,0.37,0.37}{#1}}
\newcommand{\KeywordTok}[1]{\textcolor[rgb]{0.00,0.23,0.31}{\textbf{#1}}}
\newcommand{\NormalTok}[1]{\textcolor[rgb]{0.00,0.23,0.31}{#1}}
\newcommand{\OperatorTok}[1]{\textcolor[rgb]{0.37,0.37,0.37}{#1}}
\newcommand{\OtherTok}[1]{\textcolor[rgb]{0.00,0.23,0.31}{#1}}
\newcommand{\PreprocessorTok}[1]{\textcolor[rgb]{0.68,0.00,0.00}{#1}}
\newcommand{\RegionMarkerTok}[1]{\textcolor[rgb]{0.00,0.23,0.31}{#1}}
\newcommand{\SpecialCharTok}[1]{\textcolor[rgb]{0.37,0.37,0.37}{#1}}
\newcommand{\SpecialStringTok}[1]{\textcolor[rgb]{0.13,0.47,0.30}{#1}}
\newcommand{\StringTok}[1]{\textcolor[rgb]{0.13,0.47,0.30}{#1}}
\newcommand{\VariableTok}[1]{\textcolor[rgb]{0.07,0.07,0.07}{#1}}
\newcommand{\VerbatimStringTok}[1]{\textcolor[rgb]{0.13,0.47,0.30}{#1}}
\newcommand{\WarningTok}[1]{\textcolor[rgb]{0.37,0.37,0.37}{\textit{#1}}}

\usepackage{longtable,booktabs,array}
\newcounter{none} % for unnumbered tables
\usepackage{calc} % for calculating minipage widths
% Correct order of tables after \paragraph or \subparagraph
\usepackage{etoolbox}
\makeatletter
\patchcmd\longtable{\par}{\if@noskipsec\mbox{}\fi\par}{}{}
\makeatother
% Allow footnotes in longtable head/foot
\IfFileExists{footnotehyper.sty}{\usepackage{footnotehyper}}{\usepackage{footnote}}
\makesavenoteenv{longtable}
\usepackage{graphicx}
\makeatletter
\newsavebox\pandoc@box
\newcommand*\pandocbounded[1]{% scales image to fit in text height/width
  \sbox\pandoc@box{#1}%
  \Gscale@div\@tempa{\textheight}{\dimexpr\ht\pandoc@box+\dp\pandoc@box\relax}%
  \Gscale@div\@tempb{\linewidth}{\wd\pandoc@box}%
  \ifdim\@tempb\p@<\@tempa\p@\let\@tempa\@tempb\fi% select the smaller of both
  \ifdim\@tempa\p@<\p@\scalebox{\@tempa}{\usebox\pandoc@box}%
  \else\usebox{\pandoc@box}%
  \fi%
}
% Set default figure placement to htbp
\def\fps@figure{htbp}
\makeatother



\ifLuaTeX
\usepackage[bidi=basic,shorthands=off,provide=*]{babel}
\else
\usepackage[bidi=default,shorthands=off,provide=*]{babel}
\fi
\ifPDFTeX
\else
\babelfont{rm}[]{TeX Gyre Termes}
\fi
\ifLuaTeX
  \usepackage{selnolig} % disable illegal ligatures
\fi


\setlength{\emergencystretch}{3em} % prevent overfull lines

\providecommand{\tightlist}{%
  \setlength{\itemsep}{0pt}\setlength{\parskip}{0pt}}



 


% ============================================================
% meta.tex  --  SATU-SATUNYA tempat mengisi data diri & skripsi.
% Dipakai oleh halaman judul, pernyataan, pengesahan, intisari, abstract.
% ============================================================
\usepackage{graphicx}

\newcommand{\judul}{ANALISIS PENURUNAN TEKANAN DAN TRANSFER KALOR HEAT PIPE DI REAKTOR ANTARIKSA MENGGUNAKAN PYTHON}
\newcommand{\judulinggris}{Thesis Title in English}

\newcommand{\nama}{Muhammad Naufal Tamimi}
\newcommand{\nim}{20/456315/TK/50445}
\newcommand{\tahundaftar}{2020}
\newcommand{\tahun}{2026}

% Teknik Nuklir / Teknik Fisika
\newcommand{\prodi}{Teknik Nuklir}
% Nuclear Engineering / Engineering Physics
\newcommand{\prodien}{Nuclear Engineering}

\newcommand{\kotatanggal}{Yogyakarta, tanggal-bulan-tahun}
\newcommand{\tglujian}{tanggal bulan tahun ujian}
\newcommand{\tglpengesahan}{DD BULAN TAHUN}
\newcommand{\tglujianinggris}{Month Date, year}

\newcommand{\pembimbingutama}{Nama Lengkap Pembimbing Utama}
\newcommand{\pembimbingpendamping}{Nama Lengkap Pembimbing Pendamping}
\newcommand{\ketuasidang}{Nama Lengkap Ketua Sidang}
\newcommand{\pengujiutama}{Nama Lengkap Penguji Utama}
\newcommand{\anggotapenguji}{Nama Lengkap Anggota Penguji}

% Data Ketua Departemen
\newcommand{\ketuadepartemen}{Dr. Ir. Alexander Agung, S.T., M.Sc., IPU}
\newcommand{\nipketua}{19720916 199803 1002}

% Logo UGM berdiameter 5,5 cm. Simpan sebagai fig/logo-ugm.png
\newcommand{\logougm}{%
  \IfFileExists{figures/Lambang-UGM-hitam.png}%
    {\includegraphics[width=5.5cm]{figures/Lambang-UGM-hitam.png}}%
    {\fbox{\parbox[c][5.5cm][c]{5.5cm}{\centering\small Letakkan logo UGM\\(figures/Lambang-UGM-hitam.png)}}}%
}
% ============================================================
% header.tex -- aturan tipografi mengikuti Template Skripsi 2023/2025 DTNTF UGM
% ============================================================
\usepackage{titlesec}
\usepackage{titletoc}
\usepackage{fancyhdr}
\usepackage{caption}
\usepackage{etoolbox}
\usepackage{setspace}
\usepackage{enumitem}
\usepackage{xcolor}
\usepackage{array}
\usepackage{indentfirst}

% Alinea rata kiri-kanan, awal alinea menjorok 1,27 cm (sama dengan templat)
\setlength{\parindent}{1.27cm}
\setlength{\parskip}{0pt}

% ---------- Penomoran ----------
% Bab: "BAB I"; subbab: I.1, I.2.1 (Romawi.Arab)
% Gambar/Tabel/Persamaan: Arab per bab (5.1); di lampiran: A.1
\renewcommand{\thechapter}{BAB~\Roman{chapter}}
\renewcommand{\thesection}{\Roman{chapter}.\arabic{section}}
\renewcommand{\thesubsection}{\thesection.\arabic{subsection}}
\renewcommand{\thesubsubsection}{\thesubsection.\arabic{subsubsection}}
\renewcommand{\thefigure}{\arabic{chapter}.\arabic{figure}}
\renewcommand{\thetable}{\arabic{chapter}.\arabic{table}}
\renewcommand{\theequation}{\arabic{chapter}.\arabic{equation}}

% ---------- Judul bab / subbab ----------
% Judul bab: "BAB I" lalu judul, HURUF KAPITAL, tebal, rata tengah
\titleformat{\chapter}[display]
  {\normalfont\bfseries\centering\setstretch{1}}{\thechapter}{4pt}{\MakeUppercase}
\titleformat{name=\chapter,numberless}[block]
  {\normalfont\bfseries\centering\setstretch{1}}{}{0pt}{\MakeUppercase}
\titlespacing*{\chapter}{0pt}{0pt}{22pt}

% Subbab: "I.1. Judul", tebal, rata kiri, Title Case (diketik manual)
\titleformat{\section}{\normalfont\bfseries\setstretch{1}}{\thesection.}{0.6em}{}
\titleformat{\subsection}{\normalfont\bfseries\setstretch{1}}{\thesubsection.}{0.6em}{}
\titleformat{\subsubsection}{\normalfont\bfseries\setstretch{1}}{\thesubsubsection.}{0.6em}{}
\titlespacing*{\section}{0pt}{16pt}{8pt}
\titlespacing*{\subsection}{0pt}{14pt}{8pt}
\titlespacing*{\subsubsection}{0pt}{12pt}{8pt}

% ---------- Daftar isi ----------
\titlecontents{chapter}[0pt]{\addvspace{4pt}}
  {\thecontentslabel\ \ }{}{\titlerule*[.5pc]{.}\contentspage}
\titlecontents{section}[1cm]{}
  {\thecontentslabel.\ \ }{}{\titlerule*[.5pc]{.}\contentspage}
\titlecontents{subsection}[2cm]{}
  {\thecontentslabel.\ \ }{}{\titlerule*[.5pc]{.}\contentspage}
% Tingkat khusus untuk LAMPIRAN A, B, ... (menjorok di bawah entri "LAMPIRAN")
\titlecontents{appchapter}[1cm]{}
  {\thecontentslabel\ \ }{}{\titlerule*[.5pc]{.}\contentspage}

\newcommand{\newaddcontentsline}[3]{%
  \ifstrequal{#1}{toc}%
    {\ifstrequal{#2}{chapter}%
      {\origaddcontentsline{toc}{appchapter}{#3}}%
      {\origaddcontentsline{#1}{#2}{#3}}}%
    {\origaddcontentsline{#1}{#2}{#3}}}
\newcommand{\appendixtoc}{%
  \let\origaddcontentsline\addcontentsline
  \let\addcontentsline\newaddcontentsline}

% ---------- Nomor halaman ----------
% Tengah bawah di semua halaman (sesuai templat)
\pagestyle{fancy}
\fancyhf{}
\fancyfoot[C]{\thepage}
\renewcommand{\headrulewidth}{0pt}
\fancypagestyle{plain}{\fancyhf{}\fancyfoot[C]{\thepage}\renewcommand{\headrulewidth}{0pt}}

% ---------- Judul gambar/tabel ----------
% Gambar: judul di bawah, "Gambar 5.1." label tebal, baris lanjutan menggantung.
% Tabel: judul di atas; satu baris -> rata tengah, lebih dari satu baris -> rata kiri-kanan.
\captionsetup{font={stretch=1}, labelfont=bf, labelsep=period,
              singlelinecheck=off, justification=justified, format=hang}
\captionsetup[table]{singlelinecheck=on}
\DeclareCaptionListFormat{gambarlist}{\figurename~#1#2.}
\DeclareCaptionListFormat{tabellist}{\tablename~#1#2.}
\captionsetup[figure]{listformat=gambarlist}
\captionsetup[table]{listformat=tabellist}
\makeatletter
\renewcommand*{\l@figure}{\@dottedtocline{1}{0em}{6em}}
\renewcommand*{\l@table}{\@dottedtocline{1}{0em}{6em}}
\makeatother

% ---------- Daftar (rincian): nomor menjorok 0,6 cm, teks 1,27 cm ----------
\setlist{leftmargin=1.27cm, labelsep=0.2cm, align=right, topsep=2pt}

% ---------- Daftar pustaka: spasi 1 ----------
\AtBeginEnvironment{CSLReferences}{\singlespacing}

\makeatletter
\@ifpackageloaded{caption}{}{\usepackage{caption}}
\AtBeginDocument{%
\ifdefined\contentsname
  \renewcommand*\contentsname{Daftar Isi}
\else
  \newcommand\contentsname{Daftar Isi}
\fi
\ifdefined\listfigurename
  \renewcommand*\listfigurename{Daftar Gambar}
\else
  \newcommand\listfigurename{Daftar Gambar}
\fi
\ifdefined\listtablename
  \renewcommand*\listtablename{Daftar Tabel}
\else
  \newcommand\listtablename{Daftar Tabel}
\fi
\ifdefined\figurename
  \renewcommand*\figurename{Gambar}
\else
  \newcommand\figurename{Gambar}
\fi
\ifdefined\tablename
  \renewcommand*\tablename{Tabel}
\else
  \newcommand\tablename{Tabel}
\fi
}
\@ifpackageloaded{float}{}{\usepackage{float}}
\floatstyle{ruled}
\@ifundefined{c@chapter}{\newfloat{codelisting}{h}{lop}}{\newfloat{codelisting}{h}{lop}[chapter]}
\floatname{codelisting}{Daftar}
\newcommand*\listoflistings{\listof{codelisting}{Daftar Daftar}}
\makeatother
\makeatletter
\makeatother
\makeatletter
\@ifpackageloaded{caption}{}{\usepackage{caption}}
\@ifpackageloaded{subcaption}{}{\usepackage{subcaption}}
\makeatother
\usepackage{bookmark}
\IfFileExists{xurl.sty}{\usepackage{xurl}}{} % add URL line breaks if available
\urlstyle{same}
% fallback for those not using the hyperref driver hyperxmp:
\makeatletter
\@ifundefined{xmpquote}{\newcommand{\xmpquote}[1]{#1}}{}
\makeatother
\hypersetup{
  pdflang={id},
  hidelinks,
  pdfcreator={LaTeX via pandoc}}


\author{}
\date{}
\begin{document}


\setstretch{1.43}
% Bagian Awal memakai angka Romawi kecil (i, ii, ...). Halaman ini = i (nomor tidak dicetak).
\pagenumbering{roman}
\thispagestyle{empty}
\phantomsection\addcontentsline{toc}{chapter}{HALAMAN JUDUL}
\begin{singlespace}
\begin{center}
\vspace*{0.5cm}
{\bfseries\MakeUppercase{\judul}\par}
\vspace{1.2cm}
{\bfseries SKRIPSI\par}
\vspace{1cm}
untuk memenuhi sebagian persyaratan\\
untuk memperoleh derajat Sarjana\\
Program Studi \prodi
\vspace{1.2cm}

\logougm

\vspace{1.2cm}
Diajukan oleh\\
\MakeUppercase{\nama}\\
\nim
\vspace{1.2cm}

Kepada\\
{\bfseries DEPARTEMEN TEKNIK NUKLIR DAN TEKNIK FISIKA}\\
{\bfseries FAKULTAS TEKNIK}\\
{\bfseries UNIVERSITAS GADJAH MADA}\\
{\bfseries YOGYAKARTA}\\
{\bfseries \tahun}
\end{center}
\end{singlespace}
\clearpage

\phantomsection\addcontentsline{toc}{chapter}{PERNYATAAN BEBAS PLAGIARISME}
\begin{center}
{\bfseries PERNYATAAN BEBAS PLAGIASI}
\end{center}
\vspace{0.3cm}
\begin{spacing}{1.15}
Saya yang bertanda tangan di bawah ini:

\medskip
\begin{tabular}{@{}l@{\hspace{1.6cm}}l@{\hspace{0.4cm}}l}
Nama            & : & \nama \\[3pt]
NIM             & : & \nim \\[3pt]
Tahun terdaftar & : & \tahundaftar \\[3pt]
Program Studi   & : & \prodi \\[3pt]
Fakultas        & : & Teknik \\
\end{tabular}

\medskip
menyatakan bahwa dokumen ilmiah skripsi ini tidak terdapat bagian dari karya
ilmiah lain yang telah diajukan untuk memperoleh gelar akademik di suatu lembaga
Pendidikan Tinggi, dan juga tidak terdapat karya atau pendapat yang pernah ditulis
atau diterbitkan oleh orang/lembaga lain, kecuali yang secara tertulis disitasi dalam
dokumen ini dan disebutkan sumbernya secara lengkap dalam daftar pustaka.

Dengan demikian saya menyatakan bahwa dokumen ilmiah ini bebas dari unsur-unsur
plagiasi dan apabila dokumen ilmiah Skripsi ini di kemudian hari terbukti merupakan
plagiasi dari hasil karya penulis lain dan/atau dengan sengaja mengajukan karya atau
pendapat yang merupakan hasil karya penulis lain, maka penulis bersedia menerima
sanksi akademik dan/atau sanksi hukum yang berlaku.
\end{spacing}

\vspace{0.8cm}
\begin{flushright}
\begin{minipage}{6.5cm}
\kotatanggal\\[0.3cm]
\fcolorbox{red}{red!8}{\parbox[c][2.6cm][c]{5.9cm}{\small Materai ASLI (e-Meterai)\\[2pt] Rp. 10 000\\[2pt] dan tanda tangan}}\\[0.3cm]
\nama\\
NIM. \nim
\end{minipage}
\end{flushright}
\clearpage

\phantomsection\addcontentsline{toc}{chapter}{HALAMAN PENGESAHAN}
\begin{singlespace}
\begin{center}
{\bfseries HALAMAN PENGESAHAN}\\[0.4cm]
SKRIPSI\\[0.4cm]
{\bfseries\MakeUppercase{\judul}}
\end{center}
\vspace{0.6cm}

\renewcommand{\arraystretch}{1.9}
\noindent
\begin{tabular}{@{}p{4.6cm}@{}l@{\hspace{0.3cm}}l}
Nama Mahasiswa         & : & \nama \\
Nomor Mahasiswa        & : & \nim \\
Pembimbing Utama       & : & \pembimbingutama \\
Pembimbing Pendamping  & : & \pembimbingpendamping \\
\end{tabular}

\vspace{0.5cm}
\begin{center}
Skripsi ini telah dipertahankan di depan Tim Penguji\\
pada tanggal \textit{\tglujian}
\end{center}

\noindent
\begin{tabular}{@{}p{4.6cm}@{}l@{\hspace{0.3cm}}l}
Ketua Sidang    & : & \ketuasidang \\
Penguji Utama   & : & \pengujiutama \\
Anggota Penguji & : & \anggotapenguji \\
\end{tabular}
\renewcommand{\arraystretch}{1}

\vspace{0.5cm}
\begin{center}
Skripsi ini telah diterima dan dinyatakan memenuhi\\
syarat kelulusan pada tanggal \tglpengesahan

\vspace{0.8cm}
Ketua Departemen Teknik Nuklir dan Teknik Fisika\\
Fakultas Teknik UGM\\[2.2cm]
\underline{\ketuadepartemen}\\
NIP. \nipketua
\end{center}
\end{singlespace}
\clearpage

% Halaman persembahan & motto -- ada di templat 2023/2025. Hapus berkas ini dan baris
% include-nya di skripsi.qmd bila tidak dipakai.
\phantomsection\addcontentsline{toc}{chapter}{HALAMAN PERSEMBAHAN}
\vspace*{1cm}
\begin{center}
\itshape \ldots{} Karya ini kupersembahkan untuk \ldots{} \ldots{}
\end{center}
\clearpage

\phantomsection\addcontentsline{toc}{chapter}{MOTTO}
\vspace*{1cm}
\begin{center}
\itshape \ldots{} Motto hari ini adalah \ldots{} \ldots{}
\end{center}
\clearpage

\chapter*{KATA PENGANTAR}\label{kata-pengantar}
\addcontentsline{toc}{chapter}{KATA PENGANTAR}

Kata Pengantar berisikan uraian yang mengantar pada para pembaca skripsi
kepada permasalahan yang diteliti. Dalam Kata Pengantar dapat pula
disertakan ucapan terima kasih dan apresiasi penulis kepada pihak-pihak
yang telah membantu dalam menyelesaikan tugas akhirnya. Ucapan terima
kasih disampaikan secara singkat dan harus diungkapkan dengan serius
dalam tata bahasa yang benar. Gunakan kata ``penulis'', bukan ``saya''.

\vspace{1cm}
\begin{flushright}
\kotatanggal\\[1.2cm]
Penulis
\end{flushright}
\clearpage

% Urutan: Daftar Isi, Daftar Tabel, Daftar Gambar (lalu Daftar Lambang & Singkatan).
\begin{singlespace}
\renewcommand{\contentsname}{DAFTAR ISI}
\tableofcontents
\clearpage

\phantomsection\addcontentsline{toc}{chapter}{DAFTAR TABEL}
\renewcommand{\listtablename}{DAFTAR TABEL}
\listoftables
\clearpage

\phantomsection\addcontentsline{toc}{chapter}{DAFTAR GAMBAR}
\renewcommand{\listfigurename}{DAFTAR GAMBAR}
\listoffigures
\clearpage
\end{singlespace}

\chapter*{DAFTAR LAMBANG DAN
SINGKATAN}\label{daftar-lambang-dan-singkatan}
\addcontentsline{toc}{chapter}{DAFTAR LAMBANG DAN SINGKATAN}

\textbf{Lambang Romawi}

{\def\LTcaptype{none} % do not increment counter
\begin{longtable}[]{@{}clc@{}}
\toprule\noalign{}
Lambang & Kuantitas & Satuan \\
\midrule\noalign{}
\endhead
\bottomrule\noalign{}
\endlastfoot
\(A\) & Luas penampang & m² \\
\(D_v\) & Diameter ruang uap & m \\
\(f\) & Faktor gesek Darcy & -- \\
\(g_{\mathrm{eff}}\) & Percepatan efektif & m/s² \\
\(h_{fg}\) & Kalor laten penguapan & J/kg \\
\(K\) & Permeabilitas sumbu kapiler & m² \\
\(L\) & Panjang pipa kalor & m \\
\(\dot m\) & Laju aliran massa & kg/s \\
\(P\) & Tekanan & Pa \\
\(Q\) & Laju transportasi kalor & W \\
\(r_c\) & Jari-jari pori efektif & m \\
\(R_c\) & Resistansi termal kondenser & K/W \\
\(T\) & Temperatur & K \\
\(u\) & Kecepatan uap & m/s \\
\(x\) & Koordinat aksial & m \\
\end{longtable}
}

\textbf{Lambang Yunani}

{\def\LTcaptype{none} % do not increment counter
\begin{longtable}[]{@{}clc@{}}
\toprule\noalign{}
Lambang & Kuantitas & Satuan \\
\midrule\noalign{}
\endhead
\bottomrule\noalign{}
\endlastfoot
\(\theta\) & Sudut kontak atau kemiringan, sesuai konteks & ° \\
\(\mu\) & Viskositas dinamis & Pa·s \\
\(\rho\) & Densitas & kg/m³ \\
\(\sigma\) & Tegangan permukaan & N/m \\
\end{longtable}
}

\textbf{Subskrip dan Singkatan}

{\def\LTcaptype{none} % do not increment counter
\begin{longtable}[]{@{}cl@{}}
\toprule\noalign{}
Lambang & Deskripsi \\
\midrule\noalign{}
\endhead
\bottomrule\noalign{}
\endlastfoot
\(c\) & Kondenser atau kapiler, sesuai konteks \\
\(e\) & Evaporator \\
\(l\) & Fasa cair \\
\(v\) & Fasa uap \\
CFD & \emph{Computational Fluid Dynamics} \\
HPCR & \emph{Heat Pipe Cooled Reactor} \\
NIM & Nomor Induk Mahasiswa \\
UGM & Universitas Gadjah Mada \\
\end{longtable}
}

\clearpage

\phantomsection\addcontentsline{toc}{chapter}{INTISARI}
\thispagestyle{plain}
\begin{singlespace}
\begin{center}
{\bfseries\MakeUppercase{\judul}}\\[0.4cm]
\nama\\[0.2cm]
\nim\\[0.5cm]
Diajukan kepada Departemen Teknik Nuklir dan Teknik Fisika Fakultas Teknik\\
Universitas Gadjah Mada pada tanggal \tglujian\\
untuk memenuhi sebagian persyaratan untuk memperoleh derajat\\
Sarjana Program Studi \prodi\\[0.6cm]
{\bfseries INTISARI}
\end{center}
\setlength{\parindent}{1.27cm}

\textbf{DESIGN AND ANALYSIS OF A PASSIVE HEAT-PIPE COOLING RETROFIT FOR
THE INDONESIAN SPACE REACTOR (ISR MARK 2)}

\textbf{Nama Lengkap Mahasiswa}\\
\textbf{xx/yyyyyy/TK/zzzzz}

Diajukan kepada Departemen Teknik Nuklir dan Teknik Fisika, Fakultas
Teknik, Universitas Gadjah Mada, untuk memenuhi sebagian persyaratan
penelitian program Sarjana Teknik Nuklir.

Reaktor Antariksa Indonesia (ISR) berdaya 500 kWth menggunakan pendingin
Na-78K yang dipompa aktif dan memiliki 61 pin bahan bakar dengan kanal
pendingin berdiameter 0,6 cm. Studi ini mengembangkan simulator pipa
kalor satu dimensi berbasis jaringan resistansi termal dan batas
operasional analitik, memvalidasinya terhadap data uji KRUSTY, lalu
mengevaluasi retrofit pendingin pipa kalor pasif untuk ISR Mark 2.
Simulator diimplementasikan dalam Julia dan diuji pada rentang 550--800
\(^\circ\)C; kesalahan maksimum terhadap data throughput KRUSTY adalah
12,5 \%, dengan titik operasi 295 W pada 601,5 \(^\circ\)C dan beda suhu
evaporator--kondensor 15,8 K. Untuk kanal 0,6 cm, batas pipa kalor
terbaik hanya 231 W per pin, sehingga 8,2 kW per pin tidak dapat
dicapai. Diperlukan diameter pipa 2,3--6,8 cm, bergantung pada fluida
dan suhu. Opsi A dengan natrium pada 1.200 K dan kanal 2,5 cm mencapai
9,13 kW per pin, tetapi menghilangkan 56,4 \% volume bahan bakar dan
menurunkan estimasi umur menjadi 12,2 tahun. Opsi B memerlukan kisi
teras baru dengan pitch 5,07 cm, sedangkan Opsi C hanya layak sebagai
demonstrator pasif berdaya rendah.

\textbf{Kata kunci:} pipa kalor, reaktor antariksa, KRUSTY, ISR,
natrium, pendinginan pasif

\vspace{0.3cm}
\noindent\textbf{\textit{Kata kunci:}} \textit{kata kunci 1, kata kunci 2, kata kunci 3, kata kunci 4}

\bigskip
\noindent
\begin{tabular}{@{}l@{\hspace{0.3cm}}l@{\hspace{0.3cm}}l}
Pembimbing Utama      & : & \pembimbingutama \\[4pt]
Pembimbing Pendamping & : & \pembimbingpendamping \\
\end{tabular}
\end{singlespace}
\clearpage

\phantomsection\addcontentsline{toc}{chapter}{ABSTRACT}
\thispagestyle{plain}
\begin{singlespace}
\begin{center}
{\bfseries\MakeUppercase{\judulinggris}}\\[0.4cm]
\nama\\[0.2cm]
\nim\\[0.5cm]
Submitted to the Department of Nuclear Engineering and Engineering Physics\\
Faculty of Engineering Universitas Gadjah Mada on \textit{\tglujianinggris}\\
in partial fulfillment of the requirement for the Degree of\\
Bachelor of Engineering in \prodien\\[0.6cm]
{\bfseries ABSTRACT}
\end{center}
\selectlanguage{english}

Teks dalam bahasa Inggris sangat disarankan dicek tata bahasanya
misalnya dengan grammarly.com dan cek similarity misalnya dengan
aplikasi TURNITIN.

First paragraph: background, briefly.

Second paragraph: methodology, briefly.

Third paragraph: results and conclusion, briefly.

\selectlanguage{indonesian}
\vspace{0.3cm}
\noindent\textbf{\textit{Keywords:}} \textit{Keyword 1, keyword 2, keyword 3, keyword 4}

\bigskip
\noindent
\begin{tabular}{@{}l@{\hspace{0.3cm}}l@{\hspace{0.3cm}}l}
Supervisor    & : & \pembimbingutama \\[4pt]
Co-supervisor & : & \pembimbingpendamping \\
\end{tabular}
\end{singlespace}
\clearpage

\clearpage
\pagenumbering{arabic}
\setcounter{page}{1}

\chapter{PENDAHULUAN}\label{pendahuluan}

\section{Latar Belakang}\label{latar-belakang}

Latar belakang penelitian\ldots{}

\section{Perumusan Masalah}\label{perumusan-masalah}

Akar permasalahan\ldots{}

\chapter*{DAFTAR PUSTAKA}\label{daftar-pustaka}
\addcontentsline{toc}{chapter}{DAFTAR PUSTAKA}

\protect\phantomsection\label{refs}

\chapter{LISTING PROGRAM SIMULATOR}\label{listing-program-simulator}

Lampiran ini memuat sumber lengkap \texttt{src/heatpipe\_simulator.py}.
Program menggunakan satuan SI dan menghasilkan ulang seluruh data serta
gambar yang dibahas pada Bab V.

\begin{codelisting}

\caption{\texttt{src/heatpipe\_simulator.py}}

\begin{Shaded}
\begin{Highlighting}[]
\CommentTok{\#!/usr/bin/env python3}
\CommentTok{"""Reduced{-}order, steady 1{-}D sodium heat{-}pipe simulator.}

\CommentTok{The model couples axial vapor flow, Darcy liquid return through a homogeneous}
\CommentTok{wick, and capillary pressure. It is intended for transparent thesis case}
\CommentTok{studies, not design certification or transient/CFD analysis.}

\CommentTok{Sodium correlations follow Fink and Leibowitz (1995), ANL/RE{-}95/2.}
\CommentTok{All internal calculations use SI units.}
\CommentTok{"""}
\ImportTok{from}\NormalTok{ \_\_future\_\_ }\ImportTok{import}\NormalTok{ annotations}

\ImportTok{import}\NormalTok{ argparse}
\ImportTok{import}\NormalTok{ csv}
\ImportTok{import}\NormalTok{ json}
\ImportTok{import}\NormalTok{ math}
\ImportTok{import}\NormalTok{ warnings}
\ImportTok{from}\NormalTok{ dataclasses }\ImportTok{import}\NormalTok{ asdict, dataclass, replace}
\ImportTok{from}\NormalTok{ pathlib }\ImportTok{import}\NormalTok{ Path}
\ImportTok{from}\NormalTok{ typing }\ImportTok{import}\NormalTok{ Iterable}

\ImportTok{import}\NormalTok{ matplotlib}
\NormalTok{matplotlib.use(}\StringTok{"Agg"}\NormalTok{)}
\ImportTok{import}\NormalTok{ matplotlib.pyplot }\ImportTok{as}\NormalTok{ plt}
\ImportTok{import}\NormalTok{ numpy }\ImportTok{as}\NormalTok{ np}

\NormalTok{R\_UNIVERSAL }\OperatorTok{=} \FloatTok{8.314462618}  \CommentTok{\# J mol{-}1 K{-}1}
\NormalTok{MOLAR\_MASS\_NA }\OperatorTok{=} \FloatTok{22.98976928e{-}3}  \CommentTok{\# kg mol{-}1}
\NormalTok{K\_BOLTZMANN }\OperatorTok{=} \FloatTok{1.380649e{-}23}  \CommentTok{\# J K{-}1}


\AttributeTok{@dataclass}\NormalTok{(frozen}\OperatorTok{=}\VariableTok{True}\NormalTok{)}
\KeywordTok{class}\NormalTok{ SodiumProperties:}
    \CommentTok{"""Saturated{-}sodium properties; recommended thesis range: 800{-}{-}1400 K."""}

\NormalTok{    temperature: }\BuiltInTok{float}

    \KeywordTok{def}\NormalTok{ \_\_post\_init\_\_(}\VariableTok{self}\NormalTok{) }\OperatorTok{{-}\textgreater{}} \VariableTok{None}\NormalTok{:}
        \ControlFlowTok{if} \KeywordTok{not} \FloatTok{371.0} \OperatorTok{\textless{}=} \VariableTok{self}\NormalTok{.temperature }\OperatorTok{\textless{}} \FloatTok{2503.7}\NormalTok{:}
            \ControlFlowTok{raise} \PreprocessorTok{ValueError}\NormalTok{(}\StringTok{"Temperature must be in [371, 2503.7) K."}\NormalTok{)}
        \ControlFlowTok{if} \KeywordTok{not} \FloatTok{800.0} \OperatorTok{\textless{}=} \VariableTok{self}\NormalTok{.temperature }\OperatorTok{\textless{}=} \FloatTok{1400.0}\NormalTok{:}
\NormalTok{            warnings.warn(}
                \SpecialStringTok{f"T=}\SpecialCharTok{\{}\VariableTok{self}\SpecialCharTok{.}\NormalTok{temperature}\SpecialCharTok{:.1f\}}\SpecialStringTok{ K is outside the thesis validation "}
                \StringTok{"target of 800{-}{-}1400 K; check extrapolation."}\NormalTok{,}
\NormalTok{                stacklevel}\OperatorTok{=}\DecValTok{2}\NormalTok{,}
\NormalTok{            )}

    \AttributeTok{@property}
    \KeywordTok{def}\NormalTok{ pressure(}\VariableTok{self}\NormalTok{) }\OperatorTok{{-}\textgreater{}} \BuiltInTok{float}\NormalTok{:}
        \CommentTok{"""Saturation pressure (Pa), Browning{-}{-}Potter correlation."""}
\NormalTok{        ln\_p\_mpa }\OperatorTok{=}\NormalTok{ (}
            \FloatTok{11.9463} \OperatorTok{{-}} \FloatTok{12633.73} \OperatorTok{/} \VariableTok{self}\NormalTok{.temperature}
            \OperatorTok{{-}} \FloatTok{0.4672} \OperatorTok{*}\NormalTok{ math.log(}\VariableTok{self}\NormalTok{.temperature)}
\NormalTok{        )}
        \ControlFlowTok{return}\NormalTok{ math.exp(ln\_p\_mpa) }\OperatorTok{*} \FloatTok{1.0e6}

    \AttributeTok{@property}
    \KeywordTok{def}\NormalTok{ latent\_heat(}\VariableTok{self}\NormalTok{) }\OperatorTok{{-}\textgreater{}} \BuiltInTok{float}\NormalTok{:}
        \CommentTok{"""Specific enthalpy of vaporization (J kg{-}1)."""}
\NormalTok{        tau }\OperatorTok{=} \FloatTok{1.0} \OperatorTok{{-}} \VariableTok{self}\NormalTok{.temperature }\OperatorTok{/} \FloatTok{2503.7}
\NormalTok{        h\_fg\_kj\_kg }\OperatorTok{=} \FloatTok{393.37} \OperatorTok{*}\NormalTok{ tau }\OperatorTok{+} \FloatTok{4398.6} \OperatorTok{*}\NormalTok{ tau}\OperatorTok{**}\FloatTok{0.29302}
        \ControlFlowTok{return}\NormalTok{ h\_fg\_kj\_kg }\OperatorTok{*} \FloatTok{1000.0}

    \AttributeTok{@property}
    \KeywordTok{def}\NormalTok{ liquid\_density(}\VariableTok{self}\NormalTok{) }\OperatorTok{{-}\textgreater{}} \BuiltInTok{float}\NormalTok{:}
        \CommentTok{"""Saturated liquid density (kg m{-}3)."""}
\NormalTok{        tau }\OperatorTok{=} \FloatTok{1.0} \OperatorTok{{-}} \VariableTok{self}\NormalTok{.temperature }\OperatorTok{/} \FloatTok{2503.7}
        \ControlFlowTok{return} \FloatTok{219.0} \OperatorTok{+} \FloatTok{275.32} \OperatorTok{*}\NormalTok{ tau }\OperatorTok{+} \FloatTok{511.58} \OperatorTok{*}\NormalTok{ math.sqrt(tau)}

    \AttributeTok{@property}
    \KeywordTok{def}\NormalTok{ vapor\_density(}\VariableTok{self}\NormalTok{) }\OperatorTok{{-}\textgreater{}} \BuiltInTok{float}\NormalTok{:}
        \CommentTok{"""Ideal{-}gas approximation to saturated vapor density (kg m{-}3)."""}
        \ControlFlowTok{return}\NormalTok{ (}
            \VariableTok{self}\NormalTok{.pressure }\OperatorTok{*}\NormalTok{ MOLAR\_MASS\_NA}
            \OperatorTok{/}\NormalTok{ (R\_UNIVERSAL }\OperatorTok{*} \VariableTok{self}\NormalTok{.temperature)}
\NormalTok{        )}

    \AttributeTok{@property}
    \KeywordTok{def}\NormalTok{ liquid\_viscosity(}\VariableTok{self}\NormalTok{) }\OperatorTok{{-}\textgreater{}} \BuiltInTok{float}\NormalTok{:}
        \CommentTok{"""Dynamic viscosity of liquid sodium (Pa s)."""}
        \ControlFlowTok{return}\NormalTok{ math.exp(}
            \OperatorTok{{-}}\FloatTok{6.4406} \OperatorTok{{-}} \FloatTok{0.3958} \OperatorTok{*}\NormalTok{ math.log(}\VariableTok{self}\NormalTok{.temperature)}
            \OperatorTok{+} \FloatTok{556.835} \OperatorTok{/} \VariableTok{self}\NormalTok{.temperature}
\NormalTok{        )}

    \AttributeTok{@property}
    \KeywordTok{def}\NormalTok{ vapor\_viscosity(}\VariableTok{self}\NormalTok{) }\OperatorTok{{-}\textgreater{}} \BuiltInTok{float}\NormalTok{:}
        \CommentTok{"""Hard{-}sphere dilute{-}gas estimate of vapor viscosity (Pa s).}

\CommentTok{        A 4.0 Angstrom collision diameter is used because the recommended}
\CommentTok{        liquid{-}property report does not provide a sodium{-}vapor viscosity}
\CommentTok{        correlation. Replace this approximation for high{-}fidelity studies.}
\CommentTok{        """}
\NormalTok{        mass }\OperatorTok{=}\NormalTok{ MOLAR\_MASS\_NA }\OperatorTok{/} \FloatTok{6.02214076e23}
\NormalTok{        diameter }\OperatorTok{=} \FloatTok{4.0e{-}10}
        \ControlFlowTok{return}\NormalTok{ (}
            \FloatTok{5.0} \OperatorTok{/}\NormalTok{ (}\FloatTok{16.0} \OperatorTok{*}\NormalTok{ diameter}\OperatorTok{**}\DecValTok{2}\NormalTok{)}
            \OperatorTok{*}\NormalTok{ math.sqrt(mass }\OperatorTok{*}\NormalTok{ K\_BOLTZMANN }\OperatorTok{*} \VariableTok{self}\NormalTok{.temperature }\OperatorTok{/}\NormalTok{ math.pi)}
\NormalTok{        )}

    \AttributeTok{@property}
    \KeywordTok{def}\NormalTok{ surface\_tension(}\VariableTok{self}\NormalTok{) }\OperatorTok{{-}\textgreater{}} \BuiltInTok{float}\NormalTok{:}
        \CommentTok{"""Liquid surface tension (N m{-}1), Van der Waals correlation."""}
        \ControlFlowTok{return}\NormalTok{ (}
            \FloatTok{240.5e{-}3} \OperatorTok{*}\NormalTok{ (}\FloatTok{1.0} \OperatorTok{{-}} \VariableTok{self}\NormalTok{.temperature }\OperatorTok{/} \FloatTok{2503.7}\NormalTok{)}\OperatorTok{**}\FloatTok{1.126}
\NormalTok{        )}

    \AttributeTok{@property}
    \KeywordTok{def}\NormalTok{ liquid\_conductivity(}\VariableTok{self}\NormalTok{) }\OperatorTok{{-}\textgreater{}} \BuiltInTok{float}\NormalTok{:}
        \CommentTok{"""Liquid thermal conductivity (W m{-}1 K{-}1)."""}
\NormalTok{        t }\OperatorTok{=} \VariableTok{self}\NormalTok{.temperature}
        \ControlFlowTok{return}\NormalTok{ (}
            \FloatTok{124.67} \OperatorTok{{-}} \FloatTok{0.11381} \OperatorTok{*}\NormalTok{ t }\OperatorTok{+} \FloatTok{5.5226e{-}5} \OperatorTok{*}\NormalTok{ t}\OperatorTok{**}\DecValTok{2}
            \OperatorTok{{-}} \FloatTok{1.1842e{-}8} \OperatorTok{*}\NormalTok{ t}\OperatorTok{**}\DecValTok{3}
\NormalTok{        )}

    \AttributeTok{@property}
    \KeywordTok{def}\NormalTok{ specific\_gas\_constant(}\VariableTok{self}\NormalTok{) }\OperatorTok{{-}\textgreater{}} \BuiltInTok{float}\NormalTok{:}
        \ControlFlowTok{return}\NormalTok{ R\_UNIVERSAL }\OperatorTok{/}\NormalTok{ MOLAR\_MASS\_NA}

    \AttributeTok{@property}
    \KeywordTok{def}\NormalTok{ sound\_speed(}\VariableTok{self}\NormalTok{) }\OperatorTok{{-}\textgreater{}} \BuiltInTok{float}\NormalTok{:}
        \CommentTok{"""Monatomic ideal{-}gas sound{-}speed estimate (m s{-}1)."""}
        \ControlFlowTok{return}\NormalTok{ math.sqrt((}\FloatTok{5.0} \OperatorTok{/} \FloatTok{3.0}\NormalTok{) }\OperatorTok{*} \VariableTok{self}\NormalTok{.specific\_gas\_constant }\OperatorTok{*} \VariableTok{self}\NormalTok{.temperature)}

    \KeywordTok{def}\NormalTok{ summary(}\VariableTok{self}\NormalTok{) }\OperatorTok{{-}\textgreater{}} \BuiltInTok{dict}\NormalTok{[}\BuiltInTok{str}\NormalTok{, }\BuiltInTok{float}\NormalTok{]:}
        \ControlFlowTok{return}\NormalTok{ \{}
            \StringTok{"temperature\_K"}\NormalTok{: }\VariableTok{self}\NormalTok{.temperature,}
            \StringTok{"saturation\_pressure\_Pa"}\NormalTok{: }\VariableTok{self}\NormalTok{.pressure,}
            \StringTok{"latent\_heat\_J\_kg"}\NormalTok{: }\VariableTok{self}\NormalTok{.latent\_heat,}
            \StringTok{"liquid\_density\_kg\_m3"}\NormalTok{: }\VariableTok{self}\NormalTok{.liquid\_density,}
            \StringTok{"vapor\_density\_kg\_m3"}\NormalTok{: }\VariableTok{self}\NormalTok{.vapor\_density,}
            \StringTok{"liquid\_viscosity\_Pa\_s"}\NormalTok{: }\VariableTok{self}\NormalTok{.liquid\_viscosity,}
            \StringTok{"vapor\_viscosity\_Pa\_s"}\NormalTok{: }\VariableTok{self}\NormalTok{.vapor\_viscosity,}
            \StringTok{"surface\_tension\_N\_m"}\NormalTok{: }\VariableTok{self}\NormalTok{.surface\_tension,}
            \StringTok{"liquid\_conductivity\_W\_mK"}\NormalTok{: }\VariableTok{self}\NormalTok{.liquid\_conductivity,}
            \StringTok{"sound\_speed\_m\_s"}\NormalTok{: }\VariableTok{self}\NormalTok{.sound\_speed,}
\NormalTok{        \}}


\AttributeTok{@dataclass}\NormalTok{(frozen}\OperatorTok{=}\VariableTok{True}\NormalTok{)}
\KeywordTok{class}\NormalTok{ HeatPipeGeometry:}
\NormalTok{    evaporator\_length: }\BuiltInTok{float} \OperatorTok{=} \FloatTok{0.25}
\NormalTok{    adiabatic\_length: }\BuiltInTok{float} \OperatorTok{=} \FloatTok{0.15}
\NormalTok{    condenser\_length: }\BuiltInTok{float} \OperatorTok{=} \FloatTok{0.60}
\NormalTok{    vapor\_diameter: }\BuiltInTok{float} \OperatorTok{=} \FloatTok{0.012}
\NormalTok{    wick\_thickness: }\BuiltInTok{float} \OperatorTok{=} \FloatTok{0.001}
\NormalTok{    pore\_radius: }\BuiltInTok{float} \OperatorTok{=} \FloatTok{75.0e{-}6}
\NormalTok{    permeability: }\BuiltInTok{float} \OperatorTok{=} \FloatTok{5.0e{-}10}
\NormalTok{    contact\_angle\_deg: }\BuiltInTok{float} \OperatorTok{=} \FloatTok{0.0}
\NormalTok{    effective\_wick\_conductivity: }\BuiltInTok{float} \OperatorTok{=} \FloatTok{10.0}
\NormalTok{    bubble\_radius: }\BuiltInTok{float} \OperatorTok{=} \FloatTok{5.0e{-}6}

    \KeywordTok{def}\NormalTok{ \_\_post\_init\_\_(}\VariableTok{self}\NormalTok{) }\OperatorTok{{-}\textgreater{}} \VariableTok{None}\NormalTok{:}
\NormalTok{        positive }\OperatorTok{=}\NormalTok{ \{}
\NormalTok{            key: value }\ControlFlowTok{for}\NormalTok{ key, value }\KeywordTok{in}\NormalTok{ asdict(}\VariableTok{self}\NormalTok{).items()}
            \ControlFlowTok{if}\NormalTok{ key }\OperatorTok{!=} \StringTok{"contact\_angle\_deg"}
\NormalTok{        \}}
        \ControlFlowTok{if} \BuiltInTok{any}\NormalTok{(value }\OperatorTok{\textless{}=} \FloatTok{0.0} \ControlFlowTok{for}\NormalTok{ value }\KeywordTok{in}\NormalTok{ positive.values()):}
            \ControlFlowTok{raise} \PreprocessorTok{ValueError}\NormalTok{(}\StringTok{"All geometry and wick properties must be positive."}\NormalTok{)}
        \ControlFlowTok{if} \KeywordTok{not} \FloatTok{0.0} \OperatorTok{\textless{}=} \VariableTok{self}\NormalTok{.contact\_angle\_deg }\OperatorTok{\textless{}} \FloatTok{90.0}\NormalTok{:}
            \ControlFlowTok{raise} \PreprocessorTok{ValueError}\NormalTok{(}\StringTok{"Contact angle must be in [0, 90) degrees."}\NormalTok{)}

    \AttributeTok{@property}
    \KeywordTok{def}\NormalTok{ total\_length(}\VariableTok{self}\NormalTok{) }\OperatorTok{{-}\textgreater{}} \BuiltInTok{float}\NormalTok{:}
        \ControlFlowTok{return} \VariableTok{self}\NormalTok{.evaporator\_length }\OperatorTok{+} \VariableTok{self}\NormalTok{.adiabatic\_length }\OperatorTok{+} \VariableTok{self}\NormalTok{.condenser\_length}

    \AttributeTok{@property}
    \KeywordTok{def}\NormalTok{ effective\_length(}\VariableTok{self}\NormalTok{) }\OperatorTok{{-}\textgreater{}} \BuiltInTok{float}\NormalTok{:}
        \ControlFlowTok{return}\NormalTok{ (}
            \VariableTok{self}\NormalTok{.evaporator\_length }\OperatorTok{/} \FloatTok{2.0} \OperatorTok{+} \VariableTok{self}\NormalTok{.adiabatic\_length}
            \OperatorTok{+} \VariableTok{self}\NormalTok{.condenser\_length }\OperatorTok{/} \FloatTok{2.0}
\NormalTok{        )}

    \AttributeTok{@property}
    \KeywordTok{def}\NormalTok{ vapor\_area(}\VariableTok{self}\NormalTok{) }\OperatorTok{{-}\textgreater{}} \BuiltInTok{float}\NormalTok{:}
        \ControlFlowTok{return}\NormalTok{ math.pi }\OperatorTok{*} \VariableTok{self}\NormalTok{.vapor\_diameter}\OperatorTok{**}\DecValTok{2} \OperatorTok{/} \FloatTok{4.0}

    \AttributeTok{@property}
    \KeywordTok{def}\NormalTok{ wick\_area(}\VariableTok{self}\NormalTok{) }\OperatorTok{{-}\textgreater{}} \BuiltInTok{float}\NormalTok{:}
\NormalTok{        r\_v }\OperatorTok{=} \VariableTok{self}\NormalTok{.vapor\_diameter }\OperatorTok{/} \FloatTok{2.0}
        \ControlFlowTok{return}\NormalTok{ math.pi }\OperatorTok{*}\NormalTok{ ((r\_v }\OperatorTok{+} \VariableTok{self}\NormalTok{.wick\_thickness)}\OperatorTok{**}\DecValTok{2} \OperatorTok{{-}}\NormalTok{ r\_v}\OperatorTok{**}\DecValTok{2}\NormalTok{)}

    \AttributeTok{@property}
    \KeywordTok{def}\NormalTok{ capillary\_pressure(}\VariableTok{self}\NormalTok{) }\OperatorTok{{-}\textgreater{}} \BuiltInTok{float}\NormalTok{:}
        \ControlFlowTok{raise} \PreprocessorTok{AttributeError}\NormalTok{(}\StringTok{"Capillary pressure also requires fluid temperature."}\NormalTok{)}


\AttributeTok{@dataclass}\NormalTok{(frozen}\OperatorTok{=}\VariableTok{True}\NormalTok{)}
\KeywordTok{class}\NormalTok{ OperatingConditions:}
\NormalTok{    heat\_input: }\BuiltInTok{float} \OperatorTok{=} \FloatTok{800.0}
\NormalTok{    sink\_temperature: }\BuiltInTok{float} \OperatorTok{=} \FloatTok{700.0}
\NormalTok{    condenser\_resistance: }\BuiltInTok{float} \OperatorTok{=} \FloatTok{0.375}
\NormalTok{    gravity: }\BuiltInTok{float} \OperatorTok{=} \FloatTok{0.0}
\NormalTok{    tilt\_deg: }\BuiltInTok{float} \OperatorTok{=} \FloatTok{0.0}
\NormalTok{    cells: }\BuiltInTok{int} \OperatorTok{=} \DecValTok{400}

    \KeywordTok{def}\NormalTok{ \_\_post\_init\_\_(}\VariableTok{self}\NormalTok{) }\OperatorTok{{-}\textgreater{}} \VariableTok{None}\NormalTok{:}
        \ControlFlowTok{if} \VariableTok{self}\NormalTok{.heat\_input }\OperatorTok{\textless{}} \FloatTok{0.0} \KeywordTok{or} \VariableTok{self}\NormalTok{.sink\_temperature }\OperatorTok{\textless{}=} \FloatTok{0.0}\NormalTok{:}
            \ControlFlowTok{raise} \PreprocessorTok{ValueError}\NormalTok{(}
                \StringTok{"Heat input must be nonnegative and sink temperature positive."}
\NormalTok{            )}
        \ControlFlowTok{if} \VariableTok{self}\NormalTok{.condenser\_resistance }\OperatorTok{\textless{}=} \FloatTok{0.0} \KeywordTok{or} \VariableTok{self}\NormalTok{.gravity }\OperatorTok{\textless{}} \FloatTok{0.0}\NormalTok{:}
            \ControlFlowTok{raise} \PreprocessorTok{ValueError}\NormalTok{(}
                \StringTok{"Condenser resistance must be positive and gravity nonnegative."}
\NormalTok{            )}
        \ControlFlowTok{if} \VariableTok{self}\NormalTok{.cells }\OperatorTok{\textless{}} \DecValTok{20}\NormalTok{:}
            \ControlFlowTok{raise} \PreprocessorTok{ValueError}\NormalTok{(}\StringTok{"Use at least 20 axial cells."}\NormalTok{)}

    \AttributeTok{@property}
    \KeywordTok{def}\NormalTok{ operating\_temperature(}\VariableTok{self}\NormalTok{) }\OperatorTok{{-}\textgreater{}} \BuiltInTok{float}\NormalTok{:}
        \ControlFlowTok{return} \VariableTok{self}\NormalTok{.sink\_temperature }\OperatorTok{+} \VariableTok{self}\NormalTok{.condenser\_resistance }\OperatorTok{*} \VariableTok{self}\NormalTok{.heat\_input}


\AttributeTok{@dataclass}
\KeywordTok{class}\NormalTok{ SimulationResult:}
\NormalTok{    x: np.ndarray}
\NormalTok{    section: np.ndarray}
\NormalTok{    axial\_heat\_flow: np.ndarray}
\NormalTok{    mass\_flow: np.ndarray}
\NormalTok{    vapor\_velocity: np.ndarray}
\NormalTok{    reynolds: np.ndarray}
\NormalTok{    mach: np.ndarray}
\NormalTok{    vapor\_pressure\_drop: np.ndarray}
\NormalTok{    liquid\_pressure\_drop: np.ndarray}
\NormalTok{    summary: }\BuiltInTok{dict}\NormalTok{[}\BuiltInTok{str}\NormalTok{, }\BuiltInTok{float} \OperatorTok{|} \BuiltInTok{str} \OperatorTok{|} \BuiltInTok{bool}\NormalTok{]}


\KeywordTok{def}\NormalTok{ capillary\_pressure(geometry: HeatPipeGeometry, fluid: SodiumProperties) }\OperatorTok{{-}\textgreater{}} \BuiltInTok{float}\NormalTok{:}
    \ControlFlowTok{return}\NormalTok{ (}
        \FloatTok{2.0} \OperatorTok{*}\NormalTok{ fluid.surface\_tension}
        \OperatorTok{*}\NormalTok{ math.cos(math.radians(geometry.contact\_angle\_deg))}
        \OperatorTok{/}\NormalTok{ geometry.pore\_radius}
\NormalTok{    )}


\KeywordTok{def}\NormalTok{ axial\_heat\_profile(}
\NormalTok{    x: np.ndarray, heat\_input: }\BuiltInTok{float}\NormalTok{, geometry: HeatPipeGeometry,}
\NormalTok{) }\OperatorTok{{-}\textgreater{}}\NormalTok{ np.ndarray:}
    \CommentTok{"""Heat carried by vapor; x=0 is the closed evaporator end."""}
\NormalTok{    le, la, lc }\OperatorTok{=}\NormalTok{ (}
\NormalTok{        geometry.evaporator\_length,}
\NormalTok{        geometry.adiabatic\_length,}
\NormalTok{        geometry.condenser\_length,}
\NormalTok{    )}
\NormalTok{    q }\OperatorTok{=}\NormalTok{ np.empty\_like(x)}
\NormalTok{    evap }\OperatorTok{=}\NormalTok{ x }\OperatorTok{\textless{}=}\NormalTok{ le}
\NormalTok{    adiab }\OperatorTok{=}\NormalTok{ (x }\OperatorTok{\textgreater{}}\NormalTok{ le) }\OperatorTok{\&}\NormalTok{ (x }\OperatorTok{\textless{}=}\NormalTok{ le }\OperatorTok{+}\NormalTok{ la)}
\NormalTok{    cond }\OperatorTok{=}\NormalTok{ x }\OperatorTok{\textgreater{}}\NormalTok{ le }\OperatorTok{+}\NormalTok{ la}
\NormalTok{    q[evap] }\OperatorTok{=}\NormalTok{ heat\_input }\OperatorTok{*}\NormalTok{ x[evap] }\OperatorTok{/}\NormalTok{ le}
\NormalTok{    q[adiab] }\OperatorTok{=}\NormalTok{ heat\_input}
\NormalTok{    q[cond] }\OperatorTok{=}\NormalTok{ heat\_input }\OperatorTok{*}\NormalTok{ (}\FloatTok{1.0} \OperatorTok{{-}}\NormalTok{ (x[cond] }\OperatorTok{{-}}\NormalTok{ le }\OperatorTok{{-}}\NormalTok{ la) }\OperatorTok{/}\NormalTok{ lc)}
    \ControlFlowTok{return}\NormalTok{ np.maximum(q, }\FloatTok{0.0}\NormalTok{)}


\KeywordTok{def}\NormalTok{ \_cumulative\_trapezoid(y: np.ndarray, x: np.ndarray) }\OperatorTok{{-}\textgreater{}}\NormalTok{ np.ndarray:}
\NormalTok{    result }\OperatorTok{=}\NormalTok{ np.zeros\_like(y)}
\NormalTok{    result[}\DecValTok{1}\NormalTok{:] }\OperatorTok{=}\NormalTok{ np.cumsum(}\FloatTok{0.5} \OperatorTok{*}\NormalTok{ (y[}\DecValTok{1}\NormalTok{:] }\OperatorTok{+}\NormalTok{ y[:}\OperatorTok{{-}}\DecValTok{1}\NormalTok{]) }\OperatorTok{*}\NormalTok{ np.diff(x))}
    \ControlFlowTok{return}\NormalTok{ result}


\KeywordTok{def}\NormalTok{ simulate\_at\_temperature(}
\NormalTok{    geometry: HeatPipeGeometry,}
\NormalTok{    conditions: OperatingConditions,}
\NormalTok{    temperature: }\BuiltInTok{float}\NormalTok{,}
\NormalTok{) }\OperatorTok{{-}\textgreater{}}\NormalTok{ SimulationResult:}
    \CommentTok{"""Solve the steady axial profile at a prescribed saturation temperature."""}
\NormalTok{    fluid }\OperatorTok{=}\NormalTok{ SodiumProperties(temperature)}
\NormalTok{    x }\OperatorTok{=}\NormalTok{ np.linspace(}\FloatTok{0.0}\NormalTok{, geometry.total\_length, conditions.cells }\OperatorTok{+} \DecValTok{1}\NormalTok{)}
\NormalTok{    q }\OperatorTok{=}\NormalTok{ axial\_heat\_profile(x, conditions.heat\_input, geometry)}
\NormalTok{    mdot }\OperatorTok{=}\NormalTok{ q }\OperatorTok{/}\NormalTok{ fluid.latent\_heat}
\NormalTok{    velocity }\OperatorTok{=}\NormalTok{ mdot }\OperatorTok{/}\NormalTok{ (fluid.vapor\_density }\OperatorTok{*}\NormalTok{ geometry.vapor\_area)}
\NormalTok{    reynolds }\OperatorTok{=}\NormalTok{ (}
\NormalTok{        fluid.vapor\_density }\OperatorTok{*}\NormalTok{ velocity }\OperatorTok{*}\NormalTok{ geometry.vapor\_diameter}
        \OperatorTok{/}\NormalTok{ fluid.vapor\_viscosity}
\NormalTok{    )}
\NormalTok{    mach }\OperatorTok{=}\NormalTok{ velocity }\OperatorTok{/}\NormalTok{ fluid.sound\_speed}
\NormalTok{    friction }\OperatorTok{=}\NormalTok{ np.where(}
\NormalTok{        reynolds }\OperatorTok{\textless{}} \FloatTok{2300.0}\NormalTok{,}
        \FloatTok{64.0} \OperatorTok{/}\NormalTok{ np.maximum(reynolds, }\FloatTok{1.0e{-}12}\NormalTok{),}
        \FloatTok{0.3164} \OperatorTok{/}\NormalTok{ np.maximum(reynolds, }\FloatTok{1.0e{-}12}\NormalTok{)}\OperatorTok{**}\FloatTok{0.25}\NormalTok{,}
\NormalTok{    )}
\NormalTok{    friction }\OperatorTok{=}\NormalTok{ np.where(reynolds }\OperatorTok{\textless{}} \FloatTok{1.0e{-}12}\NormalTok{, }\FloatTok{0.0}\NormalTok{, friction)}
\NormalTok{    dpv\_dx }\OperatorTok{=}\NormalTok{ (}
\NormalTok{        friction }\OperatorTok{/}\NormalTok{ geometry.vapor\_diameter}
        \OperatorTok{*} \FloatTok{0.5} \OperatorTok{*}\NormalTok{ fluid.vapor\_density }\OperatorTok{*}\NormalTok{ velocity}\OperatorTok{**}\DecValTok{2}
\NormalTok{    )}
\NormalTok{    dpl\_dx }\OperatorTok{=}\NormalTok{ (}
\NormalTok{        fluid.liquid\_viscosity }\OperatorTok{*}\NormalTok{ mdot}
        \OperatorTok{/}\NormalTok{ (fluid.liquid\_density }\OperatorTok{*}\NormalTok{ geometry.permeability }\OperatorTok{*}\NormalTok{ geometry.wick\_area)}
\NormalTok{    )}
\NormalTok{    dpv }\OperatorTok{=}\NormalTok{ \_cumulative\_trapezoid(dpv\_dx, x)}
\NormalTok{    dpl }\OperatorTok{=}\NormalTok{ \_cumulative\_trapezoid(dpl\_dx, x)}
\NormalTok{    dp\_gravity }\OperatorTok{=}\NormalTok{ (}
\NormalTok{        fluid.liquid\_density }\OperatorTok{*}\NormalTok{ conditions.gravity }\OperatorTok{*}\NormalTok{ geometry.total\_length}
        \OperatorTok{*}\NormalTok{ math.sin(math.radians(conditions.tilt\_deg))}
\NormalTok{    )}
\NormalTok{    p\_cap }\OperatorTok{=}\NormalTok{ capillary\_pressure(geometry, fluid)}
\NormalTok{    margin }\OperatorTok{=}\NormalTok{ p\_cap }\OperatorTok{{-}}\NormalTok{ dpv[}\OperatorTok{{-}}\DecValTok{1}\NormalTok{] }\OperatorTok{{-}}\NormalTok{ dpl[}\OperatorTok{{-}}\DecValTok{1}\NormalTok{] }\OperatorTok{{-}}\NormalTok{ dp\_gravity}

    \CommentTok{\# Engineering screening limits (Faghri/Chi{-}style reduced{-}order expressions).}
\NormalTok{    q\_sonic }\OperatorTok{=}\NormalTok{ (}
        \FloatTok{0.474} \OperatorTok{*}\NormalTok{ geometry.vapor\_area }\OperatorTok{*}\NormalTok{ fluid.latent\_heat}
        \OperatorTok{*}\NormalTok{ math.sqrt(fluid.vapor\_density }\OperatorTok{*}\NormalTok{ fluid.pressure)}
\NormalTok{    )}
\NormalTok{    q\_entrainment }\OperatorTok{=}\NormalTok{ (}
\NormalTok{        geometry.vapor\_area }\OperatorTok{*}\NormalTok{ fluid.latent\_heat}
        \OperatorTok{*}\NormalTok{ math.sqrt(}
\NormalTok{            fluid.vapor\_density }\OperatorTok{*}\NormalTok{ fluid.surface\_tension}
            \OperatorTok{/}\NormalTok{ (}\FloatTok{2.0} \OperatorTok{*}\NormalTok{ geometry.pore\_radius)}
\NormalTok{        )}
\NormalTok{    )}
\NormalTok{    r\_v }\OperatorTok{=}\NormalTok{ geometry.vapor\_diameter }\OperatorTok{/} \FloatTok{2.0}
\NormalTok{    r\_i }\OperatorTok{=}\NormalTok{ r\_v }\OperatorTok{+}\NormalTok{ geometry.wick\_thickness}
\NormalTok{    nucleation\_pressure }\OperatorTok{=} \BuiltInTok{max}\NormalTok{(}
        \FloatTok{2.0} \OperatorTok{*}\NormalTok{ fluid.surface\_tension }\OperatorTok{/}\NormalTok{ geometry.bubble\_radius}
        \OperatorTok{{-}}\NormalTok{ (dpv[}\OperatorTok{{-}}\DecValTok{1}\NormalTok{] }\OperatorTok{+}\NormalTok{ dpl[}\OperatorTok{{-}}\DecValTok{1}\NormalTok{] }\OperatorTok{+}\NormalTok{ dp\_gravity),}
        \FloatTok{0.0}\NormalTok{,}
\NormalTok{    )}
\NormalTok{    q\_boiling }\OperatorTok{=}\NormalTok{ (}
        \FloatTok{2.0} \OperatorTok{*}\NormalTok{ math.pi }\OperatorTok{*}\NormalTok{ geometry.evaporator\_length}
        \OperatorTok{*}\NormalTok{ geometry.effective\_wick\_conductivity }\OperatorTok{*}\NormalTok{ temperature}
        \OperatorTok{/}\NormalTok{ (fluid.latent\_heat }\OperatorTok{*}\NormalTok{ fluid.vapor\_density }\OperatorTok{*}\NormalTok{ math.log(r\_i }\OperatorTok{/}\NormalTok{ r\_v))}
        \OperatorTok{*}\NormalTok{ nucleation\_pressure}
\NormalTok{    )}
\NormalTok{    limits }\OperatorTok{=}\NormalTok{ \{}
        \StringTok{"capillary"}\NormalTok{: capillary\_limit\_at\_temperature(geometry, conditions, temperature),}
        \StringTok{"sonic"}\NormalTok{: q\_sonic,}
        \StringTok{"entrainment"}\NormalTok{: q\_entrainment,}
        \StringTok{"boiling"}\NormalTok{: q\_boiling,}
\NormalTok{    \}}
\NormalTok{    governing }\OperatorTok{=} \BuiltInTok{min}\NormalTok{(limits, key}\OperatorTok{=}\NormalTok{limits.get)}
\NormalTok{    utilization }\OperatorTok{=}\NormalTok{ conditions.heat\_input }\OperatorTok{/}\NormalTok{ limits[governing]}
\NormalTok{    section }\OperatorTok{=}\NormalTok{ np.select(}
\NormalTok{        [x }\OperatorTok{\textless{}=}\NormalTok{ geometry.evaporator\_length,}
\NormalTok{         x }\OperatorTok{\textless{}=}\NormalTok{ geometry.evaporator\_length }\OperatorTok{+}\NormalTok{ geometry.adiabatic\_length],}
\NormalTok{        [}\StringTok{"Evaporator"}\NormalTok{, }\StringTok{"Adiabatic"}\NormalTok{], default}\OperatorTok{=}\StringTok{"Condenser"}\NormalTok{,}
\NormalTok{    )}
\NormalTok{    summary: }\BuiltInTok{dict}\NormalTok{[}\BuiltInTok{str}\NormalTok{, }\BuiltInTok{float} \OperatorTok{|} \BuiltInTok{str} \OperatorTok{|} \BuiltInTok{bool}\NormalTok{] }\OperatorTok{=}\NormalTok{ \{}
        \StringTok{"temperature\_K"}\NormalTok{: temperature,}
        \StringTok{"heat\_input\_W"}\NormalTok{: conditions.heat\_input,}
        \StringTok{"max\_mass\_flow\_kg\_s"}\NormalTok{: }\BuiltInTok{float}\NormalTok{(np.}\BuiltInTok{max}\NormalTok{(mdot)),}
        \StringTok{"max\_vapor\_velocity\_m\_s"}\NormalTok{: }\BuiltInTok{float}\NormalTok{(np.}\BuiltInTok{max}\NormalTok{(velocity)),}
        \StringTok{"max\_reynolds"}\NormalTok{: }\BuiltInTok{float}\NormalTok{(np.}\BuiltInTok{max}\NormalTok{(reynolds)),}
        \StringTok{"max\_mach"}\NormalTok{: }\BuiltInTok{float}\NormalTok{(np.}\BuiltInTok{max}\NormalTok{(mach)),}
        \StringTok{"capillary\_pressure\_Pa"}\NormalTok{: p\_cap,}
        \StringTok{"vapor\_pressure\_drop\_Pa"}\NormalTok{: }\BuiltInTok{float}\NormalTok{(dpv[}\OperatorTok{{-}}\DecValTok{1}\NormalTok{]),}
        \StringTok{"liquid\_pressure\_drop\_Pa"}\NormalTok{: }\BuiltInTok{float}\NormalTok{(dpl[}\OperatorTok{{-}}\DecValTok{1}\NormalTok{]),}
        \StringTok{"gravity\_pressure\_drop\_Pa"}\NormalTok{: dp\_gravity,}
        \StringTok{"capillary\_margin\_Pa"}\NormalTok{: }\BuiltInTok{float}\NormalTok{(margin),}
        \StringTok{"capillary\_limit\_W"}\NormalTok{: limits[}\StringTok{"capillary"}\NormalTok{],}
        \StringTok{"sonic\_limit\_W"}\NormalTok{: q\_sonic,}
        \StringTok{"entrainment\_limit\_W"}\NormalTok{: q\_entrainment,}
        \StringTok{"boiling\_limit\_W"}\NormalTok{: q\_boiling,}
        \StringTok{"governing\_limit"}\NormalTok{: governing,}
        \StringTok{"limit\_utilization"}\NormalTok{: utilization,}
        \StringTok{"is\_capillary\_feasible"}\NormalTok{: }\BuiltInTok{bool}\NormalTok{(margin }\OperatorTok{\textgreater{}=} \FloatTok{0.0}\NormalTok{),}
        \StringTok{"is\_within\_all\_screening\_limits"}\NormalTok{: }\BuiltInTok{bool}\NormalTok{(utilization }\OperatorTok{\textless{}=} \FloatTok{1.0} \KeywordTok{and}\NormalTok{ margin }\OperatorTok{\textgreater{}=} \FloatTok{0.0}\NormalTok{),}
        \StringTok{"is\_incompressible\_screening\_ok"}\NormalTok{: }\BuiltInTok{bool}\NormalTok{(np.}\BuiltInTok{max}\NormalTok{(mach) }\OperatorTok{\textless{}} \FloatTok{0.3}\NormalTok{),}
        \StringTok{"gravity\_m\_s2"}\NormalTok{: conditions.gravity,}
        \StringTok{"tilt\_deg"}\NormalTok{: conditions.tilt\_deg,}
        \StringTok{"cells"}\NormalTok{: conditions.cells,}
\NormalTok{    \}}
    \ControlFlowTok{return}\NormalTok{ SimulationResult(}
\NormalTok{        x, section, q, mdot, velocity, reynolds, mach, dpv, dpl, summary,}
\NormalTok{    )}


\KeywordTok{def}\NormalTok{ simulate(}
\NormalTok{    geometry: HeatPipeGeometry, conditions: OperatingConditions,}
\NormalTok{) }\OperatorTok{{-}\textgreater{}}\NormalTok{ SimulationResult:}
    \ControlFlowTok{return}\NormalTok{ simulate\_at\_temperature(}
\NormalTok{        geometry, conditions, conditions.operating\_temperature,}
\NormalTok{    )}


\KeywordTok{def}\NormalTok{ capillary\_limit\_at\_temperature(}
\NormalTok{    geometry: HeatPipeGeometry,}
\NormalTok{    conditions: OperatingConditions,}
\NormalTok{    temperature: }\BuiltInTok{float}\NormalTok{,}
\NormalTok{) }\OperatorTok{{-}\textgreater{}} \BuiltInTok{float}\NormalTok{:}
    \CommentTok{"""Find Q for zero capillary margin using monotone bisection."""}
\NormalTok{    base }\OperatorTok{=}\NormalTok{ replace(conditions, heat\_input}\OperatorTok{=}\FloatTok{1.0}\NormalTok{)}

    \KeywordTok{def}\NormalTok{ margin(q: }\BuiltInTok{float}\NormalTok{) }\OperatorTok{{-}\textgreater{}} \BuiltInTok{float}\NormalTok{:}
        \CommentTok{\# Inline the pressure calculation to avoid recursive limit evaluation.}
\NormalTok{        fluid }\OperatorTok{=}\NormalTok{ SodiumProperties(temperature)}
\NormalTok{        local }\OperatorTok{=}\NormalTok{ replace(base, heat\_input}\OperatorTok{=}\NormalTok{q)}
\NormalTok{        x }\OperatorTok{=}\NormalTok{ np.linspace(}\FloatTok{0.0}\NormalTok{, geometry.total\_length, local.cells }\OperatorTok{+} \DecValTok{1}\NormalTok{)}
\NormalTok{        heat }\OperatorTok{=}\NormalTok{ axial\_heat\_profile(x, q, geometry)}
\NormalTok{        mdot }\OperatorTok{=}\NormalTok{ heat }\OperatorTok{/}\NormalTok{ fluid.latent\_heat}
\NormalTok{        velocity }\OperatorTok{=}\NormalTok{ mdot }\OperatorTok{/}\NormalTok{ (fluid.vapor\_density }\OperatorTok{*}\NormalTok{ geometry.vapor\_area)}
\NormalTok{        reynolds }\OperatorTok{=}\NormalTok{ (}
\NormalTok{            fluid.vapor\_density }\OperatorTok{*}\NormalTok{ velocity }\OperatorTok{*}\NormalTok{ geometry.vapor\_diameter}
            \OperatorTok{/}\NormalTok{ fluid.vapor\_viscosity}
\NormalTok{        )}
\NormalTok{        friction }\OperatorTok{=}\NormalTok{ np.where(}
\NormalTok{            reynolds }\OperatorTok{\textless{}} \FloatTok{2300.0}\NormalTok{,}
            \FloatTok{64.0} \OperatorTok{/}\NormalTok{ np.maximum(reynolds, }\FloatTok{1.0e{-}12}\NormalTok{),}
            \FloatTok{0.3164} \OperatorTok{/}\NormalTok{ np.maximum(reynolds, }\FloatTok{1.0e{-}12}\NormalTok{)}\OperatorTok{**}\FloatTok{0.25}\NormalTok{,}
\NormalTok{        )}
\NormalTok{        friction }\OperatorTok{=}\NormalTok{ np.where(reynolds }\OperatorTok{\textless{}} \FloatTok{1.0e{-}12}\NormalTok{, }\FloatTok{0.0}\NormalTok{, friction)}
\NormalTok{        dpv\_dx }\OperatorTok{=}\NormalTok{ (}
\NormalTok{            friction }\OperatorTok{/}\NormalTok{ geometry.vapor\_diameter}
            \OperatorTok{*} \FloatTok{0.5} \OperatorTok{*}\NormalTok{ fluid.vapor\_density }\OperatorTok{*}\NormalTok{ velocity}\OperatorTok{**}\DecValTok{2}
\NormalTok{        )}
\NormalTok{        dpl\_dx }\OperatorTok{=}\NormalTok{ (}
\NormalTok{            fluid.liquid\_viscosity }\OperatorTok{*}\NormalTok{ mdot}
            \OperatorTok{/}\NormalTok{ (fluid.liquid\_density }\OperatorTok{*}\NormalTok{ geometry.permeability }\OperatorTok{*}\NormalTok{ geometry.wick\_area)}
\NormalTok{        )}
\NormalTok{        dpv }\OperatorTok{=}\NormalTok{ np.trapezoid(dpv\_dx, x)}
\NormalTok{        dpl }\OperatorTok{=}\NormalTok{ np.trapezoid(dpl\_dx, x)}
\NormalTok{        dpg }\OperatorTok{=}\NormalTok{ (}
\NormalTok{            fluid.liquid\_density }\OperatorTok{*}\NormalTok{ local.gravity }\OperatorTok{*}\NormalTok{ geometry.total\_length}
            \OperatorTok{*}\NormalTok{ math.sin(math.radians(local.tilt\_deg))}
\NormalTok{        )}
        \ControlFlowTok{return}\NormalTok{ capillary\_pressure(geometry, fluid) }\OperatorTok{{-}}\NormalTok{ dpv }\OperatorTok{{-}}\NormalTok{ dpl }\OperatorTok{{-}}\NormalTok{ dpg}

    \ControlFlowTok{if}\NormalTok{ margin(}\FloatTok{0.0}\NormalTok{) }\OperatorTok{\textless{}=} \FloatTok{0.0}\NormalTok{:}
        \ControlFlowTok{return} \FloatTok{0.0}
\NormalTok{    low, high }\OperatorTok{=} \FloatTok{0.0}\NormalTok{, }\FloatTok{100.0}
    \ControlFlowTok{for}\NormalTok{ \_ }\KeywordTok{in} \BuiltInTok{range}\NormalTok{(}\DecValTok{30}\NormalTok{):}
        \ControlFlowTok{if}\NormalTok{ margin(high) }\OperatorTok{\textless{}} \FloatTok{0.0}\NormalTok{:}
            \ControlFlowTok{break}
\NormalTok{        high }\OperatorTok{*=} \FloatTok{2.0}
    \ControlFlowTok{else}\NormalTok{:}
        \ControlFlowTok{raise} \PreprocessorTok{RuntimeError}\NormalTok{(}\StringTok{"Could not bracket the capillary limit."}\NormalTok{)}
    \ControlFlowTok{for}\NormalTok{ \_ }\KeywordTok{in} \BuiltInTok{range}\NormalTok{(}\DecValTok{55}\NormalTok{):}
\NormalTok{        mid }\OperatorTok{=} \FloatTok{0.5} \OperatorTok{*}\NormalTok{ (low }\OperatorTok{+}\NormalTok{ high)}
        \ControlFlowTok{if}\NormalTok{ margin(mid) }\OperatorTok{\textgreater{}=} \FloatTok{0.0}\NormalTok{:}
\NormalTok{            low }\OperatorTok{=}\NormalTok{ mid}
        \ControlFlowTok{else}\NormalTok{:}
\NormalTok{            high }\OperatorTok{=}\NormalTok{ mid}
    \ControlFlowTok{return} \FloatTok{0.5} \OperatorTok{*}\NormalTok{ (low }\OperatorTok{+}\NormalTok{ high)}


\KeywordTok{def}\NormalTok{ operating\_limits(}
\NormalTok{    geometry: HeatPipeGeometry,}
\NormalTok{    conditions: OperatingConditions,}
\NormalTok{    temperatures: Iterable[}\BuiltInTok{float}\NormalTok{],}
\NormalTok{) }\OperatorTok{{-}\textgreater{}} \BuiltInTok{list}\NormalTok{[}\BuiltInTok{dict}\NormalTok{[}\BuiltInTok{str}\NormalTok{, }\BuiltInTok{float} \OperatorTok{|} \BuiltInTok{str}\NormalTok{]]:}
\NormalTok{    rows }\OperatorTok{=}\NormalTok{ []}
    \ControlFlowTok{for}\NormalTok{ temperature }\KeywordTok{in}\NormalTok{ temperatures:}
        \CommentTok{\# A small reference heat input makes boiling a property{-}based screening}
        \CommentTok{\# value; the other three limits are independent of the imposed load.}
\NormalTok{        probe }\OperatorTok{=}\NormalTok{ replace(conditions, heat\_input}\OperatorTok{=}\FloatTok{1.0}\NormalTok{)}
\NormalTok{        result }\OperatorTok{=}\NormalTok{ simulate\_at\_temperature(geometry, probe, }\BuiltInTok{float}\NormalTok{(temperature))}
\NormalTok{        limits }\OperatorTok{=}\NormalTok{ \{}
            \StringTok{"capillary"}\NormalTok{: }\BuiltInTok{float}\NormalTok{(result.summary[}\StringTok{"capillary\_limit\_W"}\NormalTok{]),}
            \StringTok{"sonic"}\NormalTok{: }\BuiltInTok{float}\NormalTok{(result.summary[}\StringTok{"sonic\_limit\_W"}\NormalTok{]),}
            \StringTok{"entrainment"}\NormalTok{: }\BuiltInTok{float}\NormalTok{(result.summary[}\StringTok{"entrainment\_limit\_W"}\NormalTok{]),}
            \StringTok{"boiling"}\NormalTok{: }\BuiltInTok{float}\NormalTok{(result.summary[}\StringTok{"boiling\_limit\_W"}\NormalTok{]),}
\NormalTok{        \}}
\NormalTok{        governing }\OperatorTok{=} \BuiltInTok{min}\NormalTok{(limits, key}\OperatorTok{=}\NormalTok{limits.get)}
\NormalTok{        rows.append(\{}
            \StringTok{"temperature\_K"}\NormalTok{: }\BuiltInTok{float}\NormalTok{(temperature),}
            \OperatorTok{**}\NormalTok{limits,}
            \StringTok{"maximum\_transport\_W"}\NormalTok{: limits[governing],}
            \StringTok{"governing\_limit"}\NormalTok{: governing,}
\NormalTok{        \})}
    \ControlFlowTok{return}\NormalTok{ rows}


\KeywordTok{def}\NormalTok{ write\_csv(}
\NormalTok{    path: Path, rows: }\BuiltInTok{list}\NormalTok{[}\BuiltInTok{dict}\NormalTok{], fieldnames: }\BuiltInTok{list}\NormalTok{[}\BuiltInTok{str}\NormalTok{] }\OperatorTok{|} \VariableTok{None} \OperatorTok{=} \VariableTok{None}\NormalTok{,}
\NormalTok{) }\OperatorTok{{-}\textgreater{}} \VariableTok{None}\NormalTok{:}
\NormalTok{    path.parent.mkdir(parents}\OperatorTok{=}\VariableTok{True}\NormalTok{, exist\_ok}\OperatorTok{=}\VariableTok{True}\NormalTok{)}
    \ControlFlowTok{if} \KeywordTok{not}\NormalTok{ rows:}
        \ControlFlowTok{raise} \PreprocessorTok{ValueError}\NormalTok{(}\StringTok{"No rows supplied."}\NormalTok{)}
    \ControlFlowTok{with}\NormalTok{ path.}\BuiltInTok{open}\NormalTok{(}\StringTok{"w"}\NormalTok{, newline}\OperatorTok{=}\StringTok{""}\NormalTok{, encoding}\OperatorTok{=}\StringTok{"utf{-}8"}\NormalTok{) }\ImportTok{as}\NormalTok{ handle:}
\NormalTok{        writer }\OperatorTok{=}\NormalTok{ csv.DictWriter(handle, fieldnames}\OperatorTok{=}\NormalTok{fieldnames }\KeywordTok{or} \BuiltInTok{list}\NormalTok{(rows[}\DecValTok{0}\NormalTok{]))}
\NormalTok{        writer.writeheader()}
\NormalTok{        writer.writerows(rows)}


\KeywordTok{def}\NormalTok{ configure\_plots() }\OperatorTok{{-}\textgreater{}} \VariableTok{None}\NormalTok{:}
\NormalTok{    plt.rcParams.update(\{}
        \StringTok{"font.size"}\NormalTok{: }\DecValTok{10}\NormalTok{,}
        \StringTok{"axes.grid"}\NormalTok{: }\VariableTok{True}\NormalTok{,}
        \StringTok{"grid.alpha"}\NormalTok{: }\FloatTok{0.25}\NormalTok{,}
        \StringTok{"figure.dpi"}\NormalTok{: }\DecValTok{150}\NormalTok{,}
        \StringTok{"savefig.dpi"}\NormalTok{: }\DecValTok{300}\NormalTok{,}
        \StringTok{"savefig.bbox"}\NormalTok{: }\StringTok{"tight"}\NormalTok{,}
\NormalTok{    \})}


\KeywordTok{def}\NormalTok{ plot\_axial(result: SimulationResult, path: Path) }\OperatorTok{{-}\textgreater{}} \VariableTok{None}\NormalTok{:}
\NormalTok{    configure\_plots()}
\NormalTok{    fig, axes }\OperatorTok{=}\NormalTok{ plt.subplots(}\DecValTok{3}\NormalTok{, }\DecValTok{1}\NormalTok{, figsize}\OperatorTok{=}\NormalTok{(}\FloatTok{7.0}\NormalTok{, }\FloatTok{8.2}\NormalTok{), sharex}\OperatorTok{=}\VariableTok{True}\NormalTok{)}
\NormalTok{    axes[}\DecValTok{0}\NormalTok{].plot(result.x, result.axial\_heat\_flow, color}\OperatorTok{=}\StringTok{"\#9b2226"}\NormalTok{)}
\NormalTok{    axes[}\DecValTok{0}\NormalTok{].set\_ylabel(}\StringTok{"Aliran kalor aksial (W)"}\NormalTok{)}
\NormalTok{    axes[}\DecValTok{1}\NormalTok{].plot(result.x, result.vapor\_velocity, color}\OperatorTok{=}\StringTok{"\#005f73"}\NormalTok{)}
\NormalTok{    axes[}\DecValTok{1}\NormalTok{].set\_ylabel(}\StringTok{"Kecepatan uap (m/s)"}\NormalTok{)}
\NormalTok{    axes[}\DecValTok{2}\NormalTok{].plot(result.x, result.vapor\_pressure\_drop, label}\OperatorTok{=}\StringTok{"Uap"}\NormalTok{, color}\OperatorTok{=}\StringTok{"\#ca6702"}\NormalTok{)}
\NormalTok{    axes[}\DecValTok{2}\NormalTok{].plot(result.x, result.liquid\_pressure\_drop, label}\OperatorTok{=}\StringTok{"Cairan"}\NormalTok{, color}\OperatorTok{=}\StringTok{"\#0a9396"}\NormalTok{)}
\NormalTok{    axes[}\DecValTok{2}\NormalTok{].set\_ylabel(}\StringTok{"Rugi tekanan kumulatif (Pa)"}\NormalTok{)}
\NormalTok{    axes[}\DecValTok{2}\NormalTok{].set\_xlabel(}\StringTok{"Posisi dari ujung evaporator, x (m)"}\NormalTok{)}
\NormalTok{    axes[}\DecValTok{2}\NormalTok{].legend()}
    \ControlFlowTok{for}\NormalTok{ ax }\KeywordTok{in}\NormalTok{ axes:}
        \ControlFlowTok{for}\NormalTok{ boundary }\KeywordTok{in}\NormalTok{ (}\FloatTok{0.25}\NormalTok{, }\FloatTok{0.40}\NormalTok{):}
\NormalTok{            ax.axvline(boundary, color}\OperatorTok{=}\StringTok{"0.45"}\NormalTok{, linestyle}\OperatorTok{=}\StringTok{"{-}{-}"}\NormalTok{, linewidth}\OperatorTok{=}\FloatTok{0.8}\NormalTok{)}
\NormalTok{    fig.tight\_layout()}
\NormalTok{    fig.savefig(path)}
\NormalTok{    plt.close(fig)}


\KeywordTok{def}\NormalTok{ plot\_limits(rows: }\BuiltInTok{list}\NormalTok{[}\BuiltInTok{dict}\NormalTok{], reference\_q: }\BuiltInTok{float}\NormalTok{, path: Path) }\OperatorTok{{-}\textgreater{}} \VariableTok{None}\NormalTok{:}
\NormalTok{    configure\_plots()}
\NormalTok{    temperature }\OperatorTok{=}\NormalTok{ [row[}\StringTok{"temperature\_K"}\NormalTok{] }\ControlFlowTok{for}\NormalTok{ row }\KeywordTok{in}\NormalTok{ rows]}
\NormalTok{    fig, axes }\OperatorTok{=}\NormalTok{ plt.subplots(}\DecValTok{2}\NormalTok{, }\DecValTok{1}\NormalTok{, figsize}\OperatorTok{=}\NormalTok{(}\FloatTok{7.2}\NormalTok{, }\FloatTok{6.4}\NormalTok{), sharex}\OperatorTok{=}\VariableTok{True}\NormalTok{)}
\NormalTok{    styles }\OperatorTok{=}\NormalTok{ \{}
        \StringTok{"capillary"}\NormalTok{: (}\StringTok{"Kapiler"}\NormalTok{, }\StringTok{"\#9b2226"}\NormalTok{),}
        \StringTok{"sonic"}\NormalTok{: (}\StringTok{"Sonik"}\NormalTok{, }\StringTok{"\#005f73"}\NormalTok{),}
        \StringTok{"entrainment"}\NormalTok{: (}\StringTok{"Entrainment"}\NormalTok{, }\StringTok{"\#ca6702"}\NormalTok{),}
\NormalTok{    \}}
    \ControlFlowTok{for}\NormalTok{ key, (label, color) }\KeywordTok{in}\NormalTok{ styles.items():}
\NormalTok{        axes[}\DecValTok{0}\NormalTok{].plot(temperature, [row[key] }\ControlFlowTok{for}\NormalTok{ row }\KeywordTok{in}\NormalTok{ rows], label}\OperatorTok{=}\NormalTok{label, color}\OperatorTok{=}\NormalTok{color)}
\NormalTok{    axes[}\DecValTok{0}\NormalTok{].axhline(}
\NormalTok{        reference\_q, color}\OperatorTok{=}\StringTok{"black"}\NormalTok{, linestyle}\OperatorTok{=}\StringTok{"{-}{-}"}\NormalTok{, label}\OperatorTok{=}\StringTok{"Beban kasus acuan"}\NormalTok{,}
\NormalTok{    )}
\NormalTok{    axes[}\DecValTok{0}\NormalTok{].set\_ylabel(}\StringTok{"Batas kalor (W)"}\NormalTok{)}
\NormalTok{    axes[}\DecValTok{0}\NormalTok{].legend(ncol}\OperatorTok{=}\DecValTok{2}\NormalTok{)}
\NormalTok{    axes[}\DecValTok{1}\NormalTok{].plot(temperature, [row[}\StringTok{"boiling"}\NormalTok{] }\ControlFlowTok{for}\NormalTok{ row }\KeywordTok{in}\NormalTok{ rows], color}\OperatorTok{=}\StringTok{"\#0a9396"}\NormalTok{)}
\NormalTok{    axes[}\DecValTok{1}\NormalTok{].set\_yscale(}\StringTok{"log"}\NormalTok{)}
\NormalTok{    axes[}\DecValTok{1}\NormalTok{].}\BuiltInTok{set}\NormalTok{(}
\NormalTok{        xlabel}\OperatorTok{=}\StringTok{"Temperatur operasi (K)"}\NormalTok{,}
\NormalTok{        ylabel}\OperatorTok{=}\StringTok{"Batas pendidihan (W)"}\NormalTok{,}
\NormalTok{    )}
\NormalTok{    fig.tight\_layout()}
\NormalTok{    fig.savefig(path)}
\NormalTok{    plt.close(fig)}


\KeywordTok{def}\NormalTok{ plot\_schematic(geometry: HeatPipeGeometry, path: Path) }\OperatorTok{{-}\textgreater{}} \VariableTok{None}\NormalTok{:}
    \CommentTok{"""Draw the modeled control volume and counter{-}current phase paths."""}
\NormalTok{    configure\_plots()}
\NormalTok{    fig, ax }\OperatorTok{=}\NormalTok{ plt.subplots(figsize}\OperatorTok{=}\NormalTok{(}\FloatTok{8.0}\NormalTok{, }\FloatTok{3.2}\NormalTok{))}
\NormalTok{    le, la, lc }\OperatorTok{=}\NormalTok{ (}
\NormalTok{        geometry.evaporator\_length, geometry.adiabatic\_length,}
\NormalTok{        geometry.condenser\_length,}
\NormalTok{    )}
\NormalTok{    length }\OperatorTok{=}\NormalTok{ geometry.total\_length}
\NormalTok{    ax.add\_patch(plt.Rectangle((}\DecValTok{0}\NormalTok{, }\FloatTok{0.25}\NormalTok{), length, }\FloatTok{1.5}\NormalTok{, fill}\OperatorTok{=}\VariableTok{False}\NormalTok{, linewidth}\OperatorTok{=}\FloatTok{2.0}\NormalTok{))}
\NormalTok{    ax.add\_patch(plt.Rectangle((}\DecValTok{0}\NormalTok{, }\FloatTok{0.25}\NormalTok{), length, }\FloatTok{0.22}\NormalTok{, color}\OperatorTok{=}\StringTok{"\#0a9396"}\NormalTok{, alpha}\OperatorTok{=}\FloatTok{0.75}\NormalTok{))}
\NormalTok{    ax.add\_patch(plt.Rectangle((}\DecValTok{0}\NormalTok{, }\FloatTok{1.53}\NormalTok{), length, }\FloatTok{0.22}\NormalTok{, color}\OperatorTok{=}\StringTok{"\#0a9396"}\NormalTok{, alpha}\OperatorTok{=}\FloatTok{0.75}\NormalTok{))}
\NormalTok{    ax.axvline(le, color}\OperatorTok{=}\StringTok{"0.35"}\NormalTok{, linestyle}\OperatorTok{=}\StringTok{"{-}{-}"}\NormalTok{)}
\NormalTok{    ax.axvline(le }\OperatorTok{+}\NormalTok{ la, color}\OperatorTok{=}\StringTok{"0.35"}\NormalTok{, linestyle}\OperatorTok{=}\StringTok{"{-}{-}"}\NormalTok{)}
\NormalTok{    ax.text(le }\OperatorTok{/} \DecValTok{2}\NormalTok{, }\FloatTok{2.02}\NormalTok{, }\StringTok{"Evaporator"}\NormalTok{, ha}\OperatorTok{=}\StringTok{"center"}\NormalTok{, weight}\OperatorTok{=}\StringTok{"bold"}\NormalTok{)}
\NormalTok{    ax.text(le }\OperatorTok{+}\NormalTok{ la }\OperatorTok{/} \DecValTok{2}\NormalTok{, }\FloatTok{2.02}\NormalTok{, }\StringTok{"Adiabatik"}\NormalTok{, ha}\OperatorTok{=}\StringTok{"center"}\NormalTok{, weight}\OperatorTok{=}\StringTok{"bold"}\NormalTok{)}
\NormalTok{    ax.text(le }\OperatorTok{+}\NormalTok{ la }\OperatorTok{+}\NormalTok{ lc }\OperatorTok{/} \DecValTok{2}\NormalTok{, }\FloatTok{2.02}\NormalTok{, }\StringTok{"Kondenser"}\NormalTok{, ha}\OperatorTok{=}\StringTok{"center"}\NormalTok{, weight}\OperatorTok{=}\StringTok{"bold"}\NormalTok{)}
\NormalTok{    ax.text(le }\OperatorTok{/} \DecValTok{2}\NormalTok{, }\FloatTok{2.22}\NormalTok{, }\StringTok{"Q masuk"}\NormalTok{, ha}\OperatorTok{=}\StringTok{"center"}\NormalTok{, color}\OperatorTok{=}\StringTok{"\#9b2226"}\NormalTok{)}
\NormalTok{    ax.text(le }\OperatorTok{+}\NormalTok{ la }\OperatorTok{+}\NormalTok{ lc }\OperatorTok{/} \DecValTok{2}\NormalTok{, }\FloatTok{2.22}\NormalTok{, }\StringTok{"Q keluar"}\NormalTok{, ha}\OperatorTok{=}\StringTok{"center"}\NormalTok{, color}\OperatorTok{=}\StringTok{"\#005f73"}\NormalTok{)}
    \ControlFlowTok{for}\NormalTok{ start }\KeywordTok{in}\NormalTok{ (}\FloatTok{0.04}\NormalTok{, }\FloatTok{0.22}\NormalTok{):}
\NormalTok{        ax.annotate(}\StringTok{""}\NormalTok{, xy}\OperatorTok{=}\NormalTok{(le }\OperatorTok{{-}} \FloatTok{0.02}\NormalTok{, start }\OperatorTok{+} \FloatTok{0.65}\NormalTok{), xytext}\OperatorTok{=}\NormalTok{(start, start }\OperatorTok{+} \FloatTok{0.65}\NormalTok{),}
\NormalTok{                    arrowprops}\OperatorTok{=}\BuiltInTok{dict}\NormalTok{(arrowstyle}\OperatorTok{=}\StringTok{"{-}\textgreater{}"}\NormalTok{, color}\OperatorTok{=}\StringTok{"\#ca6702"}\NormalTok{, lw}\OperatorTok{=}\FloatTok{1.8}\NormalTok{))}
\NormalTok{    ax.annotate(}\StringTok{""}\NormalTok{, xy}\OperatorTok{=}\NormalTok{(le }\OperatorTok{+}\NormalTok{ la }\OperatorTok{+}\NormalTok{ lc }\OperatorTok{{-}} \FloatTok{0.05}\NormalTok{, }\FloatTok{1.0}\NormalTok{), xytext}\OperatorTok{=}\NormalTok{(le }\OperatorTok{+} \FloatTok{0.02}\NormalTok{, }\FloatTok{1.0}\NormalTok{),}
\NormalTok{                arrowprops}\OperatorTok{=}\BuiltInTok{dict}\NormalTok{(arrowstyle}\OperatorTok{=}\StringTok{"{-}\textgreater{}"}\NormalTok{, color}\OperatorTok{=}\StringTok{"\#ca6702"}\NormalTok{, lw}\OperatorTok{=}\FloatTok{1.8}\NormalTok{))}
\NormalTok{    ax.annotate(}\StringTok{""}\NormalTok{, xy}\OperatorTok{=}\NormalTok{(le }\OperatorTok{+}\NormalTok{ la }\OperatorTok{+} \FloatTok{0.05}\NormalTok{, }\FloatTok{0.78}\NormalTok{), xytext}\OperatorTok{=}\NormalTok{(length }\OperatorTok{{-}} \FloatTok{0.05}\NormalTok{, }\FloatTok{0.78}\NormalTok{),}
\NormalTok{                arrowprops}\OperatorTok{=}\BuiltInTok{dict}\NormalTok{(arrowstyle}\OperatorTok{=}\StringTok{"{-}\textgreater{}"}\NormalTok{, color}\OperatorTok{=}\StringTok{"\#ca6702"}\NormalTok{, lw}\OperatorTok{=}\FloatTok{1.8}\NormalTok{))}
\NormalTok{    ax.annotate(}\StringTok{""}\NormalTok{, xy}\OperatorTok{=}\NormalTok{(}\FloatTok{0.05}\NormalTok{, }\FloatTok{0.36}\NormalTok{), xytext}\OperatorTok{=}\NormalTok{(length }\OperatorTok{{-}} \FloatTok{0.05}\NormalTok{, }\FloatTok{0.36}\NormalTok{),}
\NormalTok{                arrowprops}\OperatorTok{=}\BuiltInTok{dict}\NormalTok{(arrowstyle}\OperatorTok{=}\StringTok{"{-}\textgreater{}"}\NormalTok{, color}\OperatorTok{=}\StringTok{"white"}\NormalTok{, lw}\OperatorTok{=}\FloatTok{1.5}\NormalTok{))}
\NormalTok{    ax.annotate(}\StringTok{""}\NormalTok{, xy}\OperatorTok{=}\NormalTok{(}\FloatTok{0.05}\NormalTok{, }\FloatTok{1.64}\NormalTok{), xytext}\OperatorTok{=}\NormalTok{(length }\OperatorTok{{-}} \FloatTok{0.05}\NormalTok{, }\FloatTok{1.64}\NormalTok{),}
\NormalTok{                arrowprops}\OperatorTok{=}\BuiltInTok{dict}\NormalTok{(arrowstyle}\OperatorTok{=}\StringTok{"{-}\textgreater{}"}\NormalTok{, color}\OperatorTok{=}\StringTok{"white"}\NormalTok{, lw}\OperatorTok{=}\FloatTok{1.5}\NormalTok{))}
\NormalTok{    ax.text(length }\OperatorTok{/} \DecValTok{2}\NormalTok{, }\FloatTok{1.18}\NormalTok{, }\StringTok{"Uap jenuh"}\NormalTok{, ha}\OperatorTok{=}\StringTok{"center"}\NormalTok{, color}\OperatorTok{=}\StringTok{"\#8c510a"}\NormalTok{)}
\NormalTok{    ax.text(}
\NormalTok{        length }\OperatorTok{/} \DecValTok{2}\NormalTok{, }\FloatTok{0.05}\NormalTok{,}
        \StringTok{"Kapiler: cairan kembali melalui wick (arah berlawanan)"}\NormalTok{,}
\NormalTok{        ha}\OperatorTok{=}\StringTok{"center"}\NormalTok{, color}\OperatorTok{=}\StringTok{"\#075b5d"}\NormalTok{,}
\NormalTok{    )}
\NormalTok{    ax.text(}\DecValTok{0}\NormalTok{, }\OperatorTok{{-}}\FloatTok{0.18}\NormalTok{, }\StringTok{"x = 0"}\NormalTok{, ha}\OperatorTok{=}\StringTok{"center"}\NormalTok{)}
\NormalTok{    ax.text(length, }\OperatorTok{{-}}\FloatTok{0.18}\NormalTok{, }\StringTok{"x = L"}\NormalTok{, ha}\OperatorTok{=}\StringTok{"center"}\NormalTok{)}
\NormalTok{    ax.}\BuiltInTok{set}\NormalTok{(xlim}\OperatorTok{=}\NormalTok{(}\OperatorTok{{-}}\FloatTok{0.08}\NormalTok{, length }\OperatorTok{+} \FloatTok{0.08}\NormalTok{), ylim}\OperatorTok{=}\NormalTok{(}\OperatorTok{{-}}\FloatTok{0.3}\NormalTok{, }\FloatTok{2.45}\NormalTok{))}
\NormalTok{    ax.axis(}\StringTok{"off"}\NormalTok{)}
\NormalTok{    fig.tight\_layout()}
\NormalTok{    fig.savefig(path)}
\NormalTok{    plt.close(fig)}


\KeywordTok{def}\NormalTok{ plot\_gravity(rows: }\BuiltInTok{list}\NormalTok{[}\BuiltInTok{dict}\NormalTok{], path: Path) }\OperatorTok{{-}\textgreater{}} \VariableTok{None}\NormalTok{:}
\NormalTok{    configure\_plots()}
\NormalTok{    labels }\OperatorTok{=}\NormalTok{ [row[}\StringTok{"label"}\NormalTok{] }\ControlFlowTok{for}\NormalTok{ row }\KeywordTok{in}\NormalTok{ rows]}
\NormalTok{    values }\OperatorTok{=}\NormalTok{ [row[}\StringTok{"capillary\_limit\_W"}\NormalTok{] }\ControlFlowTok{for}\NormalTok{ row }\KeywordTok{in}\NormalTok{ rows]}
\NormalTok{    fig, ax }\OperatorTok{=}\NormalTok{ plt.subplots(figsize}\OperatorTok{=}\NormalTok{(}\FloatTok{7.2}\NormalTok{, }\FloatTok{4.2}\NormalTok{))}
\NormalTok{    bars }\OperatorTok{=}\NormalTok{ ax.bar(}
\NormalTok{        labels, values,}
\NormalTok{        color}\OperatorTok{=}\NormalTok{[}\StringTok{"\#0a9396"}\NormalTok{, }\StringTok{"\#0a9396"}\NormalTok{, }\StringTok{"\#ca6702"}\NormalTok{, }\StringTok{"\#ca6702"}\NormalTok{, }\StringTok{"\#9b2226"}\NormalTok{],}
\NormalTok{    )}
\NormalTok{    ax.set\_ylabel(}\StringTok{"Batas kapiler (W)"}\NormalTok{)}
\NormalTok{    ax.set\_xlabel(}\StringTok{"Percepatan efektif dan orientasi"}\NormalTok{)}
\NormalTok{    ax.tick\_params(axis}\OperatorTok{=}\StringTok{"x"}\NormalTok{, rotation}\OperatorTok{=}\DecValTok{20}\NormalTok{)}
\NormalTok{    ax.bar\_label(bars, fmt}\OperatorTok{=}\StringTok{"\%.0f"}\NormalTok{)}
\NormalTok{    fig.tight\_layout()}
\NormalTok{    fig.savefig(path)}
\NormalTok{    plt.close(fig)}


\KeywordTok{def}\NormalTok{ run\_analysis(project\_root: Path) }\OperatorTok{{-}\textgreater{}} \BuiltInTok{dict}\NormalTok{:}
\NormalTok{    geometry }\OperatorTok{=}\NormalTok{ HeatPipeGeometry()}
\NormalTok{    conditions }\OperatorTok{=}\NormalTok{ OperatingConditions()}
\NormalTok{    result }\OperatorTok{=}\NormalTok{ simulate(geometry, conditions)}
\NormalTok{    data\_dir }\OperatorTok{=}\NormalTok{ project\_root }\OperatorTok{/} \StringTok{"data"}
\NormalTok{    figure\_dir }\OperatorTok{=}\NormalTok{ project\_root }\OperatorTok{/} \StringTok{"figures"}
\NormalTok{    data\_dir.mkdir(parents}\OperatorTok{=}\VariableTok{True}\NormalTok{, exist\_ok}\OperatorTok{=}\VariableTok{True}\NormalTok{)}
\NormalTok{    figure\_dir.mkdir(parents}\OperatorTok{=}\VariableTok{True}\NormalTok{, exist\_ok}\OperatorTok{=}\VariableTok{True}\NormalTok{)}

\NormalTok{    axial\_rows }\OperatorTok{=}\NormalTok{ []}
    \ControlFlowTok{for}\NormalTok{ i }\KeywordTok{in} \BuiltInTok{range}\NormalTok{(}\BuiltInTok{len}\NormalTok{(result.x)):}
\NormalTok{        axial\_rows.append(\{}
            \StringTok{"x\_m"}\NormalTok{: result.x[i], }\StringTok{"section"}\NormalTok{: result.section[i],}
            \StringTok{"axial\_heat\_flow\_W"}\NormalTok{: result.axial\_heat\_flow[i],}
            \StringTok{"mass\_flow\_kg\_s"}\NormalTok{: result.mass\_flow[i],}
            \StringTok{"vapor\_velocity\_m\_s"}\NormalTok{: result.vapor\_velocity[i],}
            \StringTok{"reynolds"}\NormalTok{: result.reynolds[i], }\StringTok{"mach"}\NormalTok{: result.mach[i],}
            \StringTok{"vapor\_pressure\_drop\_Pa"}\NormalTok{: result.vapor\_pressure\_drop[i],}
            \StringTok{"liquid\_pressure\_drop\_Pa"}\NormalTok{: result.liquid\_pressure\_drop[i],}
\NormalTok{        \})}
\NormalTok{    write\_csv(data\_dir }\OperatorTok{/} \StringTok{"profil\_aksial.csv"}\NormalTok{, axial\_rows)}

\NormalTok{    temperatures }\OperatorTok{=}\NormalTok{ np.arange(}\FloatTok{800.0}\NormalTok{, }\FloatTok{1200.1}\NormalTok{, }\FloatTok{25.0}\NormalTok{)}
\NormalTok{    limit\_rows }\OperatorTok{=}\NormalTok{ operating\_limits(geometry, conditions, temperatures)}
\NormalTok{    write\_csv(data\_dir }\OperatorTok{/} \StringTok{"batas\_operasi.csv"}\NormalTok{, limit\_rows)}

\NormalTok{    load\_rows }\OperatorTok{=}\NormalTok{ []}
    \ControlFlowTok{for}\NormalTok{ heat\_input }\KeywordTok{in}\NormalTok{ (}\DecValTok{300}\NormalTok{, }\DecValTok{400}\NormalTok{, }\DecValTok{600}\NormalTok{, }\DecValTok{800}\NormalTok{, }\DecValTok{1000}\NormalTok{, }\DecValTok{1200}\NormalTok{, }\DecValTok{1400}\NormalTok{, }\DecValTok{1600}\NormalTok{):}
\NormalTok{        local }\OperatorTok{=}\NormalTok{ replace(conditions, heat\_input}\OperatorTok{=}\NormalTok{heat\_input)}
\NormalTok{        row }\OperatorTok{=}\NormalTok{ simulate(geometry, local).summary}
\NormalTok{        load\_rows.append(}\BuiltInTok{dict}\NormalTok{(row))}
\NormalTok{    write\_csv(data\_dir }\OperatorTok{/} \StringTok{"studi\_beban.csv"}\NormalTok{, load\_rows)}

\NormalTok{    gravity\_rows }\OperatorTok{=}\NormalTok{ []}
    \ControlFlowTok{for}\NormalTok{ label, gravity, tilt }\KeywordTok{in}\NormalTok{ (}
\NormalTok{        (}\StringTok{"Mikrogravitasi}\CharTok{\textbackslash{}n}\StringTok{0"}\NormalTok{, }\FloatTok{0.0}\NormalTok{, }\FloatTok{0.0}\NormalTok{),}
\NormalTok{        (}\StringTok{"Akselerasi kecil}\CharTok{\textbackslash{}n}\StringTok{0,01; 90°"}\NormalTok{, }\FloatTok{0.01}\NormalTok{, }\FloatTok{90.0}\NormalTok{),}
\NormalTok{        (}\StringTok{"Akselerasi}\CharTok{\textbackslash{}n}\StringTok{0,1; 90°"}\NormalTok{, }\FloatTok{0.1}\NormalTok{, }\FloatTok{90.0}\NormalTok{),}
\NormalTok{        (}\StringTok{"Akselerasi}\CharTok{\textbackslash{}n}\StringTok{1; 90°"}\NormalTok{, }\FloatTok{1.0}\NormalTok{, }\FloatTok{90.0}\NormalTok{),}
\NormalTok{        (}\StringTok{"Bumi}\CharTok{\textbackslash{}n}\StringTok{9,81; 90°"}\NormalTok{, }\FloatTok{9.81}\NormalTok{, }\FloatTok{90.0}\NormalTok{),}
\NormalTok{    ):}
\NormalTok{        local }\OperatorTok{=}\NormalTok{ replace(conditions, gravity}\OperatorTok{=}\NormalTok{gravity, tilt\_deg}\OperatorTok{=}\NormalTok{tilt)}
\NormalTok{        q\_cap }\OperatorTok{=}\NormalTok{ capillary\_limit\_at\_temperature(geometry, local, }\FloatTok{1000.0}\NormalTok{)}
\NormalTok{        gravity\_rows.append(\{}
            \StringTok{"label"}\NormalTok{: label, }\StringTok{"gravity\_m\_s2"}\NormalTok{: gravity, }\StringTok{"tilt\_deg"}\NormalTok{: tilt,}
            \StringTok{"capillary\_limit\_W"}\NormalTok{: q\_cap,}
\NormalTok{        \})}
\NormalTok{    write\_csv(data\_dir }\OperatorTok{/} \StringTok{"studi\_gravitasi.csv"}\NormalTok{, gravity\_rows)}

\NormalTok{    mesh\_rows }\OperatorTok{=}\NormalTok{ []}
\NormalTok{    reference\_dp }\OperatorTok{=} \VariableTok{None}
    \ControlFlowTok{for}\NormalTok{ cells }\KeywordTok{in}\NormalTok{ (}\DecValTok{50}\NormalTok{, }\DecValTok{100}\NormalTok{, }\DecValTok{200}\NormalTok{, }\DecValTok{400}\NormalTok{, }\DecValTok{800}\NormalTok{, }\DecValTok{1600}\NormalTok{):}
\NormalTok{        local }\OperatorTok{=}\NormalTok{ replace(conditions, cells}\OperatorTok{=}\NormalTok{cells)}
\NormalTok{        summary }\OperatorTok{=}\NormalTok{ simulate(geometry, local).summary}
\NormalTok{        total\_dp }\OperatorTok{=}\NormalTok{ (}
            \BuiltInTok{float}\NormalTok{(summary[}\StringTok{"vapor\_pressure\_drop\_Pa"}\NormalTok{])}
            \OperatorTok{+} \BuiltInTok{float}\NormalTok{(summary[}\StringTok{"liquid\_pressure\_drop\_Pa"}\NormalTok{])}
\NormalTok{        )}
        \ControlFlowTok{if}\NormalTok{ cells }\OperatorTok{==} \DecValTok{1600}\NormalTok{:}
\NormalTok{            reference\_dp }\OperatorTok{=}\NormalTok{ total\_dp}
\NormalTok{        mesh\_rows.append(\{}\StringTok{"cells"}\NormalTok{: cells, }\StringTok{"total\_pressure\_drop\_Pa"}\NormalTok{: total\_dp\})}
    \ControlFlowTok{for}\NormalTok{ row }\KeywordTok{in}\NormalTok{ mesh\_rows:}
\NormalTok{        row[}\StringTok{"relative\_difference\_percent"}\NormalTok{] }\OperatorTok{=}\NormalTok{ (}
            \BuiltInTok{abs}\NormalTok{(row[}\StringTok{"total\_pressure\_drop\_Pa"}\NormalTok{] }\OperatorTok{{-}}\NormalTok{ reference\_dp)}
            \OperatorTok{/}\NormalTok{ reference\_dp }\OperatorTok{*} \FloatTok{100.0}
\NormalTok{        )}
\NormalTok{    write\_csv(data\_dir }\OperatorTok{/} \StringTok{"konvergensi\_grid.csv"}\NormalTok{, mesh\_rows)}

\NormalTok{    pore\_rows }\OperatorTok{=}\NormalTok{ []}
    \ControlFlowTok{for}\NormalTok{ pore\_um }\KeywordTok{in}\NormalTok{ (}\DecValTok{25}\NormalTok{, }\DecValTok{50}\NormalTok{, }\DecValTok{75}\NormalTok{, }\DecValTok{100}\NormalTok{, }\DecValTok{150}\NormalTok{):}
\NormalTok{        local\_geometry }\OperatorTok{=}\NormalTok{ replace(geometry, pore\_radius}\OperatorTok{=}\NormalTok{pore\_um }\OperatorTok{*} \FloatTok{1.0e{-}6}\NormalTok{)}
\NormalTok{        pore\_rows.append(\{}
            \StringTok{"pore\_radius\_um"}\NormalTok{: pore\_um,}
            \StringTok{"capillary\_limit\_W"}\NormalTok{: capillary\_limit\_at\_temperature(}
\NormalTok{                local\_geometry, conditions, }\FloatTok{1000.0}\NormalTok{,}
\NormalTok{            ),}
            \StringTok{"entrainment\_limit\_W"}\NormalTok{: simulate\_at\_temperature(}
\NormalTok{                local\_geometry, replace(conditions, heat\_input}\OperatorTok{=}\FloatTok{1.0}\NormalTok{), }\FloatTok{1000.0}\NormalTok{,}
\NormalTok{            ).summary[}\StringTok{"entrainment\_limit\_W"}\NormalTok{],}
\NormalTok{        \})}
\NormalTok{    write\_csv(data\_dir }\OperatorTok{/} \StringTok{"studi\_pori.csv"}\NormalTok{, pore\_rows)}

\NormalTok{    plot\_schematic(geometry, figure\_dir }\OperatorTok{/} \StringTok{"skema{-}model.png"}\NormalTok{)}
\NormalTok{    plot\_axial(result, figure\_dir }\OperatorTok{/} \StringTok{"profil{-}aksial.png"}\NormalTok{)}
\NormalTok{    plot\_limits(limit\_rows, conditions.heat\_input, figure\_dir }\OperatorTok{/} \StringTok{"batas{-}operasi.png"}\NormalTok{)}
\NormalTok{    plot\_gravity(gravity\_rows, figure\_dir }\OperatorTok{/} \StringTok{"pengaruh{-}gravitasi.png"}\NormalTok{)}

\NormalTok{    output }\OperatorTok{=}\NormalTok{ \{}
        \StringTok{"model"}\NormalTok{: }\StringTok{"Steady 1{-}D vapor/Darcy{-}wick capillary model"}\NormalTok{,}
        \StringTok{"geometry"}\NormalTok{: asdict(geometry),}
        \StringTok{"conditions"}\NormalTok{: asdict(conditions),}
        \StringTok{"fluid"}\NormalTok{: SodiumProperties(conditions.operating\_temperature).summary(),}
        \StringTok{"results"}\NormalTok{: result.summary,}
        \StringTok{"derived\_geometry"}\NormalTok{: \{}
            \StringTok{"total\_length\_m"}\NormalTok{: geometry.total\_length,}
            \StringTok{"effective\_length\_m"}\NormalTok{: geometry.effective\_length,}
            \StringTok{"vapor\_area\_m2"}\NormalTok{: geometry.vapor\_area,}
            \StringTok{"wick\_area\_m2"}\NormalTok{: geometry.wick\_area,}
\NormalTok{        \},}
\NormalTok{    \}}
    \ControlFlowTok{with}\NormalTok{ (data\_dir }\OperatorTok{/} \StringTok{"ringkasan\_kasus.json"}\NormalTok{).}\BuiltInTok{open}\NormalTok{(}\StringTok{"w"}\NormalTok{, encoding}\OperatorTok{=}\StringTok{"utf{-}8"}\NormalTok{) }\ImportTok{as}\NormalTok{ handle:}
\NormalTok{        json.dump(output, handle, indent}\OperatorTok{=}\DecValTok{2}\NormalTok{, ensure\_ascii}\OperatorTok{=}\VariableTok{False}\NormalTok{)}
    \ControlFlowTok{return}\NormalTok{ output}


\KeywordTok{def}\NormalTok{ main() }\OperatorTok{{-}\textgreater{}} \VariableTok{None}\NormalTok{:}
\NormalTok{    parser }\OperatorTok{=}\NormalTok{ argparse.ArgumentParser(description}\OperatorTok{=}\NormalTok{\_\_doc\_\_)}
\NormalTok{    parser.add\_argument(}
        \StringTok{"{-}{-}output{-}dir"}\NormalTok{, }\BuiltInTok{type}\OperatorTok{=}\NormalTok{Path, default}\OperatorTok{=}\NormalTok{Path(}\VariableTok{\_\_file\_\_}\NormalTok{).resolve().parents[}\DecValTok{1}\NormalTok{],}
        \BuiltInTok{help}\OperatorTok{=}\StringTok{"Project directory receiving data/ and figures/ (default: project root)."}\NormalTok{,}
\NormalTok{    )}
\NormalTok{    args }\OperatorTok{=}\NormalTok{ parser.parse\_args()}
\NormalTok{    output }\OperatorTok{=}\NormalTok{ run\_analysis(args.output\_dir.resolve())}
    \BuiltInTok{print}\NormalTok{(json.dumps(output[}\StringTok{"results"}\NormalTok{], indent}\OperatorTok{=}\DecValTok{2}\NormalTok{, ensure\_ascii}\OperatorTok{=}\VariableTok{False}\NormalTok{))}


\ControlFlowTok{if} \VariableTok{\_\_name\_\_} \OperatorTok{==} \StringTok{"\_\_main\_\_"}\NormalTok{:}
\NormalTok{    main()}
\end{Highlighting}
\end{Shaded}

\end{codelisting}

\chapter{PENGUJIAN DAN PANDUAN
PENGGUNAAN}\label{pengujian-dan-panduan-penggunaan}

\section{Listing Pengujian}\label{listing-pengujian}

Berkas \texttt{tests/test\_simulator.py} berisi sepuluh pengujian
berbasis \texttt{unittest}.

\begin{codelisting}

\caption{\texttt{tests/test\_simulator.py}}

\begin{Shaded}
\begin{Highlighting}[]
\CommentTok{"""Regression and physical{-}consistency tests for the heat{-}pipe model."""}
\ImportTok{from}\NormalTok{ \_\_future\_\_ }\ImportTok{import}\NormalTok{ annotations}

\ImportTok{import}\NormalTok{ math}
\ImportTok{import}\NormalTok{ sys}
\ImportTok{import}\NormalTok{ unittest}
\ImportTok{from}\NormalTok{ dataclasses }\ImportTok{import}\NormalTok{ replace}
\ImportTok{from}\NormalTok{ pathlib }\ImportTok{import}\NormalTok{ Path}

\ImportTok{import}\NormalTok{ numpy }\ImportTok{as}\NormalTok{ np}

\NormalTok{sys.path.insert(}\DecValTok{0}\NormalTok{, }\BuiltInTok{str}\NormalTok{(Path(}\VariableTok{\_\_file\_\_}\NormalTok{).resolve().parents[}\DecValTok{1}\NormalTok{] }\OperatorTok{/} \StringTok{"src"}\NormalTok{))}

\ImportTok{from}\NormalTok{ heatpipe\_simulator }\ImportTok{import}\NormalTok{ (  }\CommentTok{\# noqa: E402}
\NormalTok{    HeatPipeGeometry,}
\NormalTok{    OperatingConditions,}
\NormalTok{    SodiumProperties,}
\NormalTok{    axial\_heat\_profile,}
\NormalTok{    capillary\_limit\_at\_temperature,}
\NormalTok{    simulate,}
\NormalTok{)}


\KeywordTok{class}\NormalTok{ TestSodiumProperties(unittest.TestCase):}
    \KeywordTok{def}\NormalTok{ test\_normal\_boiling\_point(}\VariableTok{self}\NormalTok{) }\OperatorTok{{-}\textgreater{}} \VariableTok{None}\NormalTok{:}
        \CommentTok{\# Recommended correlation gives approximately 1 atm at 1156 K.}
\NormalTok{        pressure }\OperatorTok{=}\NormalTok{ SodiumProperties(}\FloatTok{1156.0}\NormalTok{).pressure}
        \VariableTok{self}\NormalTok{.assertAlmostEqual(pressure }\OperatorTok{/} \FloatTok{101325.0}\NormalTok{, }\FloatTok{1.0}\NormalTok{, delta}\OperatorTok{=}\FloatTok{0.03}\NormalTok{)}

    \KeywordTok{def}\NormalTok{ test\_property\_values\_at\_1000\_k(}\VariableTok{self}\NormalTok{) }\OperatorTok{{-}\textgreater{}} \VariableTok{None}\NormalTok{:}
\NormalTok{        fluid }\OperatorTok{=}\NormalTok{ SodiumProperties(}\FloatTok{1000.0}\NormalTok{)}
        \VariableTok{self}\NormalTok{.assertAlmostEqual(fluid.surface\_tension, }\FloatTok{0.1354}\NormalTok{, delta}\OperatorTok{=}\FloatTok{0.0002}\NormalTok{)}
        \VariableTok{self}\NormalTok{.assertAlmostEqual(fluid.liquid\_viscosity, }\FloatTok{1.81e{-}4}\NormalTok{, delta}\OperatorTok{=}\FloatTok{2.0e{-}6}\NormalTok{)}
        \VariableTok{self}\NormalTok{.assertGreater(fluid.latent\_heat, }\FloatTok{3.8e6}\NormalTok{)}
        \VariableTok{self}\NormalTok{.assertLess(fluid.latent\_heat, }\FloatTok{4.3e6}\NormalTok{)}
        \VariableTok{self}\NormalTok{.assertGreater(fluid.liquid\_density, }\FloatTok{750.0}\NormalTok{)}
        \VariableTok{self}\NormalTok{.assertLess(fluid.liquid\_density, }\FloatTok{810.0}\NormalTok{)}

    \KeywordTok{def}\NormalTok{ test\_positive\_properties(}\VariableTok{self}\NormalTok{) }\OperatorTok{{-}\textgreater{}} \VariableTok{None}\NormalTok{:}
\NormalTok{        fluid }\OperatorTok{=}\NormalTok{ SodiumProperties(}\FloatTok{1000.0}\NormalTok{)}
        \ControlFlowTok{for}\NormalTok{ value }\KeywordTok{in}\NormalTok{ fluid.summary().values():}
            \VariableTok{self}\NormalTok{.assertTrue(math.isfinite(value) }\KeywordTok{and}\NormalTok{ value }\OperatorTok{\textgreater{}} \FloatTok{0.0}\NormalTok{)}


\KeywordTok{class}\NormalTok{ TestHeatPipeModel(unittest.TestCase):}
    \KeywordTok{def}\NormalTok{ setUp(}\VariableTok{self}\NormalTok{) }\OperatorTok{{-}\textgreater{}} \VariableTok{None}\NormalTok{:}
        \VariableTok{self}\NormalTok{.geometry }\OperatorTok{=}\NormalTok{ HeatPipeGeometry()}
        \VariableTok{self}\NormalTok{.conditions }\OperatorTok{=}\NormalTok{ OperatingConditions()}

    \KeywordTok{def}\NormalTok{ test\_axial\_energy\_profile(}\VariableTok{self}\NormalTok{) }\OperatorTok{{-}\textgreater{}} \VariableTok{None}\NormalTok{:}
\NormalTok{        x }\OperatorTok{=}\NormalTok{ np.linspace(}\FloatTok{0.0}\NormalTok{, }\VariableTok{self}\NormalTok{.geometry.total\_length, }\DecValTok{1001}\NormalTok{)}
\NormalTok{        heat }\OperatorTok{=}\NormalTok{ axial\_heat\_profile(x, }\FloatTok{800.0}\NormalTok{, }\VariableTok{self}\NormalTok{.geometry)}
        \VariableTok{self}\NormalTok{.assertAlmostEqual(heat[}\DecValTok{0}\NormalTok{], }\FloatTok{0.0}\NormalTok{)}
        \VariableTok{self}\NormalTok{.assertAlmostEqual(heat[}\OperatorTok{{-}}\DecValTok{1}\NormalTok{], }\FloatTok{0.0}\NormalTok{)}
        \VariableTok{self}\NormalTok{.assertAlmostEqual(}\BuiltInTok{float}\NormalTok{(np.}\BuiltInTok{max}\NormalTok{(heat)), }\FloatTok{800.0}\NormalTok{)}
        \CommentTok{\# Heat is generated only in the evaporator and removed in the condenser.}
\NormalTok{        derivative }\OperatorTok{=}\NormalTok{ np.diff(heat) }\OperatorTok{/}\NormalTok{ np.diff(x)}
        \VariableTok{self}\NormalTok{.assertGreater(}\BuiltInTok{float}\NormalTok{(np.}\BuiltInTok{max}\NormalTok{(derivative[:}\DecValTok{250}\NormalTok{])), }\FloatTok{0.0}\NormalTok{)}
        \VariableTok{self}\NormalTok{.assertLess(}\BuiltInTok{float}\NormalTok{(np.}\BuiltInTok{min}\NormalTok{(derivative[}\OperatorTok{{-}}\DecValTok{600}\NormalTok{:])), }\FloatTok{0.0}\NormalTok{)}

    \KeywordTok{def}\NormalTok{ test\_microgravity\_removes\_hydrostatic\_term(}\VariableTok{self}\NormalTok{) }\OperatorTok{{-}\textgreater{}} \VariableTok{None}\NormalTok{:}
\NormalTok{        result }\OperatorTok{=}\NormalTok{ simulate(}\VariableTok{self}\NormalTok{.geometry, }\VariableTok{self}\NormalTok{.conditions)}
        \VariableTok{self}\NormalTok{.assertEqual(result.summary[}\StringTok{"gravity\_pressure\_drop\_Pa"}\NormalTok{], }\FloatTok{0.0}\NormalTok{)}
        \VariableTok{self}\NormalTok{.assertTrue(result.summary[}\StringTok{"is\_capillary\_feasible"}\NormalTok{])}

    \KeywordTok{def}\NormalTok{ test\_reference\_case\_is\_capillary\_limited(}\VariableTok{self}\NormalTok{) }\OperatorTok{{-}\textgreater{}} \VariableTok{None}\NormalTok{:}
\NormalTok{        result }\OperatorTok{=}\NormalTok{ simulate(}\VariableTok{self}\NormalTok{.geometry, }\VariableTok{self}\NormalTok{.conditions)}
        \VariableTok{self}\NormalTok{.assertEqual(result.summary[}\StringTok{"governing\_limit"}\NormalTok{], }\StringTok{"capillary"}\NormalTok{)}
        \VariableTok{self}\NormalTok{.assertLess(}\BuiltInTok{float}\NormalTok{(result.summary[}\StringTok{"limit\_utilization"}\NormalTok{]), }\FloatTok{1.0}\NormalTok{)}
        \VariableTok{self}\NormalTok{.assertLess(}\BuiltInTok{float}\NormalTok{(result.summary[}\StringTok{"max\_mach"}\NormalTok{]), }\FloatTok{0.3}\NormalTok{)}

    \KeywordTok{def}\NormalTok{ test\_capillary\_limit\_has\_zero\_margin(}\VariableTok{self}\NormalTok{) }\OperatorTok{{-}\textgreater{}} \VariableTok{None}\NormalTok{:}
\NormalTok{        q\_cap }\OperatorTok{=}\NormalTok{ capillary\_limit\_at\_temperature(}\VariableTok{self}\NormalTok{.geometry, }\VariableTok{self}\NormalTok{.conditions, }\FloatTok{1000.0}\NormalTok{)}
\NormalTok{        local }\OperatorTok{=}\NormalTok{ replace(}
            \VariableTok{self}\NormalTok{.conditions,}
\NormalTok{            heat\_input}\OperatorTok{=}\NormalTok{q\_cap,}
\NormalTok{            sink\_temperature}\OperatorTok{=}\FloatTok{1000.0} \OperatorTok{{-}} \VariableTok{self}\NormalTok{.conditions.condenser\_resistance }\OperatorTok{*}\NormalTok{ q\_cap,}
\NormalTok{        )}
\NormalTok{        result }\OperatorTok{=}\NormalTok{ simulate(}\VariableTok{self}\NormalTok{.geometry, local)}
        \VariableTok{self}\NormalTok{.assertAlmostEqual(}
            \BuiltInTok{float}\NormalTok{(result.summary[}\StringTok{"temperature\_K"}\NormalTok{]), }\FloatTok{1000.0}\NormalTok{, delta}\OperatorTok{=}\FloatTok{1e{-}8}\NormalTok{,}
\NormalTok{        )}
        \VariableTok{self}\NormalTok{.assertAlmostEqual(}
            \BuiltInTok{float}\NormalTok{(result.summary[}\StringTok{"capillary\_margin\_Pa"}\NormalTok{]) }\OperatorTok{/} \FloatTok{3612.0}\NormalTok{,}
            \FloatTok{0.0}\NormalTok{, delta}\OperatorTok{=}\FloatTok{1e{-}6}\NormalTok{,}
\NormalTok{        )}

    \KeywordTok{def}\NormalTok{ test\_adverse\_gravity\_reduces\_capacity(}\VariableTok{self}\NormalTok{) }\OperatorTok{{-}\textgreater{}} \VariableTok{None}\NormalTok{:}
\NormalTok{        q\_micro }\OperatorTok{=}\NormalTok{ capillary\_limit\_at\_temperature(}\VariableTok{self}\NormalTok{.geometry, }\VariableTok{self}\NormalTok{.conditions, }\FloatTok{1000.0}\NormalTok{)}
\NormalTok{        adverse }\OperatorTok{=}\NormalTok{ replace(}\VariableTok{self}\NormalTok{.conditions, gravity}\OperatorTok{=}\FloatTok{1.0}\NormalTok{, tilt\_deg}\OperatorTok{=}\FloatTok{90.0}\NormalTok{)}
\NormalTok{        q\_adverse }\OperatorTok{=}\NormalTok{ capillary\_limit\_at\_temperature(}\VariableTok{self}\NormalTok{.geometry, adverse, }\FloatTok{1000.0}\NormalTok{)}
        \VariableTok{self}\NormalTok{.assertLess(q\_adverse, q\_micro)}
\NormalTok{        earth }\OperatorTok{=}\NormalTok{ replace(}\VariableTok{self}\NormalTok{.conditions, gravity}\OperatorTok{=}\FloatTok{9.81}\NormalTok{, tilt\_deg}\OperatorTok{=}\FloatTok{90.0}\NormalTok{)}
        \VariableTok{self}\NormalTok{.assertEqual(}
\NormalTok{            capillary\_limit\_at\_temperature(}\VariableTok{self}\NormalTok{.geometry, earth, }\FloatTok{1000.0}\NormalTok{), }\FloatTok{0.0}\NormalTok{,}
\NormalTok{        )}

    \KeywordTok{def}\NormalTok{ test\_grid\_convergence(}\VariableTok{self}\NormalTok{) }\OperatorTok{{-}\textgreater{}} \VariableTok{None}\NormalTok{:}
\NormalTok{        coarse }\OperatorTok{=}\NormalTok{ simulate(}\VariableTok{self}\NormalTok{.geometry, replace(}\VariableTok{self}\NormalTok{.conditions, cells}\OperatorTok{=}\DecValTok{100}\NormalTok{))}
\NormalTok{        fine }\OperatorTok{=}\NormalTok{ simulate(}\VariableTok{self}\NormalTok{.geometry, replace(}\VariableTok{self}\NormalTok{.conditions, cells}\OperatorTok{=}\DecValTok{1600}\NormalTok{))}
\NormalTok{        dp\_coarse }\OperatorTok{=}\NormalTok{ (}
            \BuiltInTok{float}\NormalTok{(coarse.summary[}\StringTok{"vapor\_pressure\_drop\_Pa"}\NormalTok{])}
            \OperatorTok{+} \BuiltInTok{float}\NormalTok{(coarse.summary[}\StringTok{"liquid\_pressure\_drop\_Pa"}\NormalTok{])}
\NormalTok{        )}
\NormalTok{        dp\_fine }\OperatorTok{=}\NormalTok{ (}
            \BuiltInTok{float}\NormalTok{(fine.summary[}\StringTok{"vapor\_pressure\_drop\_Pa"}\NormalTok{])}
            \OperatorTok{+} \BuiltInTok{float}\NormalTok{(fine.summary[}\StringTok{"liquid\_pressure\_drop\_Pa"}\NormalTok{])}
\NormalTok{        )}
        \VariableTok{self}\NormalTok{.assertLess(}\BuiltInTok{abs}\NormalTok{(dp\_coarse }\OperatorTok{{-}}\NormalTok{ dp\_fine) }\OperatorTok{/}\NormalTok{ dp\_fine, }\FloatTok{0.005}\NormalTok{)}

    \KeywordTok{def}\NormalTok{ test\_invalid\_geometry\_rejected(}\VariableTok{self}\NormalTok{) }\OperatorTok{{-}\textgreater{}} \VariableTok{None}\NormalTok{:}
        \ControlFlowTok{with} \VariableTok{self}\NormalTok{.assertRaises(}\PreprocessorTok{ValueError}\NormalTok{):}
\NormalTok{            HeatPipeGeometry(permeability}\OperatorTok{=}\FloatTok{0.0}\NormalTok{)}


\ControlFlowTok{if} \VariableTok{\_\_name\_\_} \OperatorTok{==} \StringTok{"\_\_main\_\_"}\NormalTok{:}
\NormalTok{    unittest.main()}
\end{Highlighting}
\end{Shaded}

\end{codelisting}

\section{Menjalankan Program}\label{menjalankan-program}

Dari direktori proyek, gunakan perintah berikut.

\begin{Shaded}
\begin{Highlighting}[]
\ExtensionTok{python} \AttributeTok{{-}m}\NormalTok{ pip install }\AttributeTok{{-}r}\NormalTok{ requirements.txt}
\ExtensionTok{python}\NormalTok{ src/heatpipe\_simulator.py}
\ExtensionTok{python} \AttributeTok{{-}m}\NormalTok{ unittest discover }\AttributeTok{{-}s}\NormalTok{ tests }\AttributeTok{{-}v}
\end{Highlighting}
\end{Shaded}

Keluaran konsol merangkum kasus acuan dalam JSON. Berkas ringkasan kasus
pada direktori data menyimpan masukan, sifat fluida, geometri turunan,
dan hasil. Berkas CSV lain menyimpan profil aksial serta studi
parameter. Jika parameter pada dataclass diubah, seluruh data dan gambar
harus dihasilkan ulang sebelum angka dalam naskah diperbarui.

\section{Interpretasi Status Program}\label{interpretasi-status-program}

\begin{description}
\tightlist
\item[\texttt{is\_capillary\_feasible}]
Margin kapiler tidak negatif.
\item[\texttt{is\_within\_all\_screening\_limits}]
Daya berada di bawah batas pengatur dan kapiler layak menurut model.
\item[\texttt{is\_incompressible\_screening\_ok}]
Mach maksimum lebih kecil dari 0,3.
\item[\texttt{governing\_limit}]
Batas yang bernilai minimum pada kasus tersebut.
\item[\texttt{limit\_utilization}]
Rasio daya terhadap batas pengatur.
\end{description}

Status bernilai benar tidak berarti desain telah tervalidasi. Pengguna
harus memeriksa rentang sifat, asumsi geometri, kondisi batas, dan
ketidakpastian korelasi.




\end{document}
